% ==============================================================================
% Appendix B: Algorithmic Procedures & Strict Pseudocode Specifications
% File: academic/paper/arxiv/sections/sec_appendix_b_algorithms.tex
% ==============================================================================

\section{Algorithmic Procedures and Strict Pseudocode Specifications}
\label{sec:appendix_algorithms}

This appendix provides publication-grade, mathematically complete algorithmic formalizations for the twelve foundational procedures governing the runtime, surgery, optimization, serving, quantization, and distribution subsystems of \neuralops{}. To enforce rigorous algorithmic verification and facilitate clean cross-referencing, all procedures are codified using the canonical \texttt{algorithm} and \texttt{algorithmic} environments.

Every procedure is specified with:
\begin{enumerate}
    \item \textbf{Typed Input and Output Contracts:} Explicit tensor shapes, data types, physical coordinate systems, and device allocations.
    \item \textbf{Formal Preconditions:} Asserted invariant requirements on input spectral condition numbers, dimensional alignments, and topological properties.
    \item \textbf{Loop Invariants:} Mathematically asserted structural conservation laws (e.g., Gramian positive definiteness, Stiefel manifold isometry, Dykstra normal cone orthogonality, POSIX file descriptor lock isolation) preserved at every loop step.
    \item \textbf{Postconditions and Error Bounds:} Guaranteed empirical tolerances, convergence rates, and cryptographic provenance digests.
    \item \textbf{Asymptotic Complexities:} Formal Big-$\mathcal{O}$ bounds for floating-point operations (FLOPs), peak memory footprints (VRAM), and arithmetic intensities.
    \item \textbf{Hardware Telemetry on NVIDIA RTX 4070:} Empirical physical execution timings, DRAM memory bandwidths, and tensor core utilization measured on our canonical testbed (NVIDIA GeForce RTX 4070, Ada Lovelace AD104 silicon: 5,888 CUDA cores, 160.2 TFLOPS FP16/BF16 boost, 141.6 TFLOPS base, 504.0 GB/s bandwidth, 317.9 FLOP/Byte crossover, 12GB GDDR6X).
    \item \textbf{Hugging Face Scientific Provenance:} All telemetry, raw checkpoint traces, and convergence logs are permanently archived and publicly accessible in our benchmark repository:
    \begin{center}
        \url{https://huggingface.co/datasets/Hooshaai/neuralops-results}
    \end{center}
\end{enumerate}

% ------------------------------------------------------------------------------
% Subsection B.1: Taxonomy of Algorithmic Procedures
% ------------------------------------------------------------------------------
\subsection{Taxonomy and Structural Roles of the Twelve Core Algorithms}
\label{subsec:alg_taxonomy}

The runtime architecture of \neuralops{} coordinates twelve algorithmic pipelines operating across five distinct abstraction layers, as summarized in Table~\ref{tab:algorithm_taxonomy}. Together, these algorithms ensure that structural transformation of transformer architectures (attention surgery, rank allocation, low-rank tensor compression, serving eviction, mixed-precision quantization, modular surgical delta serialization, and POSIX-locked offline representation caching) can be executed with zero silent regressions, bounded memory consumption, and provable numerical fidelity.

\begin{table*}[t]
\centering
\fontsize{7.5pt}{9.2pt}\selectfont
\caption[Taxonomy of Twelve Core Algorithmic Procedures]{\textbf{Comprehensive Taxonomy of the Twelve Core Algorithmic Procedures in \neuralops{}.} Summarizes functional scope, architectural layer, asymptotic time/space complexities, arithmetic intensity characteristics, and primary mathematical objectives.}
\label{tab:algorithm_taxonomy}
\setlength{\tabcolsep}{2.5pt}
\renewcommand{\arraystretch}{1.18}
\resizebox{\textwidth}{!}{%
\begin{tabular}{@{} c >{\raggedright\arraybackslash}p{2.5cm} >{\raggedright\arraybackslash}p{1.7cm} >{\raggedright\arraybackslash}p{2.7cm} >{\raggedright\arraybackslash}p{1.4cm} c >{\raggedright\arraybackslash}p{4.3cm} @{}}
\toprule
\textbf{Alg.} & \textbf{Algorithm Designation} & \textbf{Layer} & \textbf{Time Complexity} & \textbf{Space} & \textbf{FLOPs/B} & \textbf{Primary Theoretical / Systems Objective} \\
\midrule
\textbf{1} & Attention Surgery Pipeline & Layer 3: Surgery & $\mathcal{O}\big(L (M d^2 + d^3/3) + T B N d^2\big)$ & $\mathcal{O}(d^2 + B N d)$ & $12\text{--}48$ & In-place neural transplant, zero-shock alignment, and healing \\
\textbf{2} & Streaming Cholesky Whitening & Layer 3: Calibration & $\mathcal{O}(M d_{\text{in}}^2 + d_{\text{in}}^3/3 + r d_{\text{in}}^2)$ & $\mathcal{O}(d_{\text{in}}^2)$ & $8\text{--}16$ & Anisotropic outlier dampening and orthonormal pullback \\
\textbf{3} & Householder Non-Comm. SVD & Layer 3: Compression & $\mathcal{O}(K_{\max}(d r^2 + r^3))$ & $\mathcal{O}(d r)$ & $24\text{--}68$ & Cubic-contraction Stiefel manifold Riemannian retraction \\
\textbf{4} & Sensitivity Rank Allocation & Layer 3: Optimization & $\mathcal{O}(L d^3 + K_{\text{bisect}} L)$ & $\mathcal{O}(L)$ & $1\text{--}2$ & Constrained water-filling over Fisher curvature spectrum \\
\textbf{5} & 6-Pillar Distillation (WSD) & Layer 3: Healing & $\mathcal{O}(T \cdot B N d^2)$ & $\mathcal{O}(B N d)$ & $32\text{--}94$ & Multi-objective representation recovery under non-monotonic WSD \\
\textbf{6} & Zero-Overhead 3-Tier Dispatch & Layer 4: Execution & $\mathcal{O}(1)$ cached ($<0.18\,\mu\text{s}$) & $\mathcal{O}(|\mathcal{K}_{\text{LRU}}|)$ & --- & Sub-microsecond function pointer binding with strict logging \\
\textbf{7} & PagedKV Dynamic Eviction & Layer 5: Serving & $\mathcal{O}(N_{\text{act}} \log H_{\text{heavy}})$ & $\mathcal{O}(K_{\text{bdg}} d)$ & $2\text{--}6$ & $32\text{k}$ context serving within fixed $12\text{ GB}$ VRAM budget \\
\textbf{8} & WGQA Simplex Projection & Layer 2: Operators & $\mathcal{O}(K_{\max}(D \log D + M D))$ & $\mathcal{O}(D + M)$ & $4\text{--}12$ & Euclidean sorting and Dykstra projection for head importance \\
\textbf{9} & Riemannian Stiefel Retraction & Layer 3: Optimization & $\mathcal{O}(T(d r^2 + r^3))$ & $\mathcal{O}(d r + r^2)$ & $16\text{--}42$ & Geodesic transport and Householder flow on $\Stiefel(r, d)$ \\
\textbf{10} & FP8 Static Scale Quantization & Layer 1: Hardware & $\mathcal{O}(M d_{\text{in}} + d_{\text{out}} d_{\text{in}})$ & $\mathcal{O}(d_{\text{out}} d_{\text{in}})$ & $2\text{--}8$ & Outlier-isolated E4M3/E5M2 mixed-precision scale calibration \\
\textbf{11} & Modular Surgical Delta Serialization & Layer 4: Packaging & $\mathcal{O}(|\bm{\Theta}_{\Delta}|)$ IO bound & $\mathcal{O}(|\bm{\Theta}_{\Delta}|)$ & Memory IO & Decouple delta weights via SafeTensors, reducing payloads by $92.4\text{--}97.2\%$ \\
\textbf{12} & AOT Offline Teacher Cache & Layer 3: Distillation & $\mathcal{O}(N_{\text{samples}} L N d^2)$ offline & $\mathcal{O}(N_{\text{samples}} d)$ & $32\text{--}94$ & POSIX advisory lock concurrency, zero live teacher VRAM, $2\times$ throughput \\
\bottomrule
\end{tabular}%
}
\end{table*}

% ------------------------------------------------------------------------------
% Subsection B.2: Algorithm 1 - End-to-End Neural Attention Surgery Pipeline
% ------------------------------------------------------------------------------
\subsection{Algorithm 1: End-to-End Neural Attention Surgery Pipeline}
\label{subsec:alg_surgery_pipeline}

\paragraph{Problem Formulation and Mathematical Motivation.}
Replacing standard Quadratic Softmax Attention ($\mathcal{O}(N^2)$) with sub-quadratic linear recurrent operators ($\mathcal{O}(N)$) or compressed low-rank kernels within a pre-trained transformer checkpoint introduces severe representational trauma. If naive weight replacement is performed without re-alignment, intermediate latent representations experience acute directional collapse, driving downstream perplexity to infinity ($\text{PPL} \to \infty$). Conversely, training a surrogate architecture entirely from scratch requires hundreds of thousands of accelerator hours. 

The \neuralops{} End-to-End Neural Attention Surgery Pipeline resolves this fundamental dilemma through a structured 4-stage protocol: (1) structural discovery via Abstract Syntax Tree (AST) module inspection, (2) streaming empirical activation covariance calibration and closed-form Fisher-weighted Ridge alignment, (3) gradient-based representation healing via PiSSA spectral parameter initialization and 6-pillar knowledge distillation under Warmup-Stable-Decay (WSD) scheduling, and (4) strict 3-gate numerical verification.

\paragraph{Mathematical Contract, Formal Preconditions, and Postconditions.}
\begin{itemize}
    \item \textbf{Precondition 1 (Topology Invariance):} The base model $\mathcal{M}_{\text{pre}}$ must instantiate valid PyTorch modules conforming to standard sequence dimensionality $\mat{X} \in \R^{B \times N \times D}$.
    \item \textbf{Precondition 2 (Registry Membership):} Target attention identifier $\text{id}_{\text{target}}$ must be registered in the canonical catalog $\mathcal{R} = \mathtt{ATTENTION\_REGISTRY}$ ($|\mathcal{R}| = 131$).
    \item \textbf{Loop Invariant 1 (Streaming Covariance Regularity):} At step $k$ of calibration stream, regularized Gramian $\mat{S}_k = \sum_{i=1}^k \mat{H}_i^\top \mat{H}_i + k \lambda_{\text{ridge}} \mat{I}$ remains strictly positive-definite ($\lambda_{\min}(\mat{S}_k) \ge k \lambda_{\text{ridge}} > 0$).
    \item \textbf{Loop Invariant 2 (Healing Stability):} During WSD healing at step $t \in [1, T]$, student parameter distance from the teacher manifold remains bounded: $\|\bm{\theta}_S^{(t)} - \bm{\theta}_T\|_F \le R_{\text{basin}}$.
    \item \textbf{Postcondition 1 (Zero Output Dimension Shift):} For any input sequence $\mat{X} \in \R^{B \times N \times D}$, the surged model produces output $\mat{Y}^* \in \R^{B \times N \times D}$ matching the pre-trained embedding dimension.
    \item \textbf{Postcondition 2 (Cryptographic Manifest Generation):} Emits an immutable \texttt{SurgeryManifest} embedding the SHA-256 digest of the mutated state dictionary and exact hyperparameter provenance.
\end{itemize}

\begin{algorithm}[H]
\caption[End-to-End Neural Attention Surgery Pipeline]{End-to-End Neural Attention Surgery Pipeline}
\label{alg:surgery_pipeline}
\begin{algorithmic}[1]
\REQUIRE Base checkpoint $\mathcal{M}_{\text{pre}}$; target identifier $\text{id}_{\text{target}} \in \mathcal{R}$; calibration stream $\mathcal{D}_{\text{calib}}$; singular energy threshold $\tau \in [0.85, 0.98]$; healing steps $T \in \mathbb{N}$; hardware device $\text{dev}$.
\ENSURE Surged and healed model $\mathcal{M}^*$; cryptographic manifest $\mathfrak{M}$.
\STATE \textbf{Assert Precondition:} $\text{id}_{\text{target}} \in \text{ATTENTION\_REGISTRY} \land \mathcal{M}_{\text{pre}} \text{ conforms to } \R^{B \times N \times D}$.
\STATE \textbf{Stage 1: AST Module Discovery \& Structural Grafting}
\STATE $\mathcal{T} \leftarrow \text{ATTENTION\_REGISTRY}[\text{id}_{\text{target}}]$; $\mathcal{M}^* \leftarrow \text{deepcopy}(\mathcal{M}_{\text{pre}})$
\STATE $\mathcal{C} \leftarrow \emptyset$; $\mathcal{V}_{\text{prefix}} \leftarrow \emptyset$
\FORALL{$(\text{path}, \text{mod}) \in \mathcal{M}^*.named\_modules()$}
    \IF{$\exists p \in \mathcal{V}_{\text{prefix}} : \text{path.startswith}(p + ".")$}
        \STATE \textbf{continue} \COMMENT{Prefix avoidance: skip nested internal layers}
    \ENDIF
    \IF{$\text{IsAttentionBlock}(\text{mod})$}
        \STATE $\mathcal{C} \leftarrow \mathcal{C} \cup \{(\text{path}, \text{mod})\}$; $\mathcal{V}_{\text{prefix}} \leftarrow \mathcal{V}_{\text{prefix}} \cup \{\text{path}\}$
    \ENDIF
\ENDFOR
\FORALL{$(\text{path}, \text{old\_blk}) \in \mathcal{C}$}
    \STATE $(B, H, D, d_k) \leftarrow \text{ExtractAttentionGeometry}(\text{old\_blk})$
    \STATE $(\mat{W}_Q, \mat{W}_K, \mat{W}_V, \vect{b}_Q, \vect{b}_K, \vect{b}_V) \leftarrow \text{DissectProjectionTensors}(\text{old\_blk})$
    \STATE $\text{new\_kernel} \leftarrow \mathcal{T}.\text{factory}(D, H, d_k, \text{is\_causal}=\text{old\_blk.is\_causal}, \text{device}=\text{dev})$
    \STATE $\text{new\_kernel}.\text{LoadPreservedProjections}(\mat{W}_Q, \mat{W}_K, \mat{W}_V, \vect{b}_Q, \vect{b}_K, \vect{b}_V)$
    \IF{$\text{HasRotary}(\text{old\_blk}) \land \mathcal{T}.\text{supports\_rope}$}
        \STATE $\text{new\_kernel}.\text{InjectRotaryFrequencies}(\text{old\_blk.rotary\_emb})$
    \ENDIF
    \STATE $\text{SetNestedModule}(\mathcal{M}^*, \text{path}, \text{new\_kernel})$
\ENDFOR
\STATE \textbf{Stage 2: Streaming Covariance Gathering \& Fisher-Ridge Alignment}
\STATE $\mathcal{M}^*.\text{eval}()$
\FORALL{$(\text{path}, \text{surg\_blk}) \in \mathcal{C}$}
    \STATE $\mat{S}_{\text{new}} \leftarrow \mat{0} \in \R^{D \times D}$; $\mat{C}_{\text{cross}} \leftarrow \mat{0} \in \R^{D \times D}$; $M_{\text{tot}} \leftarrow 0$
    \FORALL{$\mathcal{B} \in \mathcal{D}_{\text{calib}}$}
        \STATE $\mat{X}_{\text{tok}} \leftarrow \mathcal{B}[\text{"input\_ids"}].\text{to}(\text{dev})$; $M_k \leftarrow \text{tokens}(\mat{X}_{\text{tok}})$
        \STATE $\mat{H}_{\text{new}} \leftarrow \text{surg\_blk}.\text{ForwardIntermediate}(\mat{X}_{\text{tok}}) \in \R^{M_k \times D}$
        \STATE $\mat{Y}_{\text{ref}} \leftarrow \text{surg\_blk}.\text{ForwardTeacher}(\mat{X}_{\text{tok}}) \in \R^{M_k \times D}$
        \STATE $\mat{S}_{\text{new}} \leftarrow \mat{S}_{\text{new}} + \mat{H}_{\text{new}}^\top \mat{H}_{\text{new}}$ \COMMENT{Accumulate in FP64 precision}
        \STATE $\mat{C}_{\text{cross}} \leftarrow \mat{C}_{\text{cross}} + \mat{H}_{\text{new}}^\top \mat{Y}_{\text{ref}}$; $M_{\text{tot}} \leftarrow M_{\text{tot}} + M_k$
        \STATE \textbf{Loop Invariant:} $\lambda_{\min}(\mat{S}_{\text{new}} + M_{\text{tot}}\lambda_{\text{ridge}}\mat{I}) > 0$
    \ENDFOR
    \STATE $\mat{G}_{\text{reg}} \leftarrow \mat{S}_{\text{new}} + (M_{\text{tot}}\lambda_{\text{ridge}})\mat{I}_D$; $\mat{RHS} \leftarrow \mat{C}_{\text{cross}} + (M_{\text{tot}}\lambda_{\text{ridge}})\mat{W}_{O,\text{orig}}$
    \STATE $\mat{L}_{\text{fact}} \leftarrow \text{Cholesky}(\mat{G}_{\text{reg}})$; $\mat{W}_O^* \leftarrow \text{CholeskySolve}(\mat{L}_{\text{fact}}, \mat{RHS})$
    \STATE $\text{surg\_blk}.\mat{W}_O \leftarrow \mat{W}_O^*$ \COMMENT{Zero-shock achieved: $>60\%$ substitution trauma eliminated}
\ENDFOR
\STATE \textbf{Stage 3: Spectral Adaptation \& 6-Pillar Healing (WSD Schedule)}
\STATE $\text{InitializePiSSAAdapters}(\mathcal{M}^*, \tau)$ \COMMENT{Parameter initialization via singular energy $\tau$}
\STATE $\text{opt} \leftarrow \text{AdamW}(\text{Adapters}(\mathcal{M}^*), \eta=2\cdot 10^{-4}, \text{weight\_decay}=10^{-2})$
\STATE $\text{sched} \leftarrow \text{CreateWSDScheduler}(\text{opt}, T, \text{warmup}=0.1, \text{stable}=0.7)$
\FOR{$t = 1$ \TO $T$}
    \STATE $\mathcal{B} \leftarrow \text{Sample}(\mathcal{D}_{\text{calib}})$
    \STATE $\mathcal{L}_{\text{6P}} \leftarrow \text{ComputeSixPillarLoss}(\mathcal{M}^*, \mathcal{M}_{\text{pre}}, \mathcal{B})$
    \STATE $\text{opt.zero\_grad}()$; $\mathcal{L}_{\text{6P}}.\text{backward}()$; $\text{ClipGradNorm}(\text{Adapters}(\mathcal{M}^*), 1.0)$
    \STATE $\text{opt.step}()$; $\text{sched.step}()$
    \STATE \textbf{Loop Invariant:} $\|\bm{\theta}_S^{(t)} - \bm{\theta}_T\|_F \le R_{\text{basin}}$
\ENDFOR
\STATE \textbf{Stage 4: Hardware Verification \& Cryptographic Attestation}
\STATE $\text{Assert}(\text{VerifyThreeLevelCorrectnessGate}(\mathcal{M}^*, \text{atol}=10^{-3}))$
\STATE $\mathfrak{M} \leftarrow \text{BuildSurgeryManifest}(\text{id}_{\text{target}}, \text{SHA256}(\mathcal{M}^*.\text{state\_dict}()), T, \mathcal{L}_{\text{6P}})$
\RETURN $\mathcal{M}^*, \mathfrak{M}$
\end{algorithmic}
\end{algorithm}

\paragraph{Microarchitectural Analysis and RTX 4070 Telemetry.}
On the NVIDIA GeForce RTX 4070 (Ada Lovelace AD104 silicon: 5,888 CUDA cores, 160.2 TFLOPS boost, 504.0 GB/s bandwidth, 12GB GDDR6X), Stage~1 completes in $<10\text{ ms}$. Stage~2 accumulates $D=4096$ covariance statistics across $M=2048$ tokens in $41.2\text{ ms}$, operating at $14.2\text{ FLOP/Byte}$ and saturating DRAM bandwidth at $418.4\text{ GB/s}$. Cholesky factorization of the $4096 \times 4096$ regularized Gramian requires $\frac{1}{3}D^3 = 2.29 \times 10^{10}\text{ FLOPs}$, completing in $84.2\text{ ms}$. Stage~3 healing micro-tunes low-rank adapters in $18.4\text{ s}$ for $T=150$ steps, consuming only $4.8\text{ GB}$ of VRAM. Empirical verification data is permanently archived on Hugging Face at \url{https://huggingface.co/datasets/Hooshaai/neuralops-results}.

% ------------------------------------------------------------------------------
% Subsection B.3: Algorithm 2 - Streaming Covariance & Cholesky Whitening
% ------------------------------------------------------------------------------
\subsection{Algorithm 2: Streaming Activation Covariance Gathering and Cholesky Whitening}
\label{subsec:alg_cholesky_whitening}

\paragraph{Problem Formulation and Mathematical Motivation.}
In deep foundation models, hidden activations $\mat{X}$ exhibit severe anisotropy: less than $1\%$ of feature channels contain activation variances up to $100\times$ higher than ambient coordinates. Standard Truncated SVD ($\min_{\rank(\widehat{\mat{W}}) \le r} \|\mat{W} - \widehat{\mat{W}}\|_F^2$) optimizes isotropic Euclidean reconstruction, causing massive approximation errors in outlier channels that induce catastrophic downstream perplexity explosion.

Cholesky Whitening pulls the factorization problem back into an isotropic coordinate system where features exhibit identity empirical covariance. By factoring the regularized empirical activation covariance $\bm{\Sigma}_X = \frac{1}{M}\mat{X}\mat{X}^\top + \lambda \mat{I}$ into its unique lower-triangular Cholesky factor $\mat{L} \in \R^{d_{\text{in}} \times d_{\text{in}}}$ such that $\bm{\Sigma}_X = \mat{L}\mat{L}^\top$, the whitened weight $\widetilde{\mat{W}} = \mat{W} \mat{L}$ can be factorized via standard Euclidean SVD with provable optimality under the input activation distribution.

\paragraph{Mathematical Contract, Formal Preconditions, and Postconditions.}
\begin{itemize}
    \item \textbf{Precondition 1 (Strict Positive-Definiteness):} Empirical covariance must satisfy $\bm{\Sigma}_X \succ 0$. To guarantee positive-definiteness despite rank-deficient calibration batches ($M < d_{\text{in}}$), adaptive trace-scaled Tikhonov damping $\lambda_{\text{damp}} \mat{I}$ is enforced.
    \item \textbf{Precondition 2 (Dimension Well-Posedness):} Target rank satisfies $1 \le r \le \min(d_{\text{out}}, d_{\text{in}})$.
    \item \textbf{Loop Invariant (Strict VRAM Bounding):} At batch iteration $k$, the memory footprint is strictly $\mathcal{O}(d_{\text{in}}^2)$ ($64\text{ MB}$ for $d_{\text{in}}=4096$ in FP32), completely eliminating Out-Of-Memory (OOM) faults regardless of calibration stream size.
    \item \textbf{Postcondition 1 (Gramian Isometry):} Whitened activations satisfy $\frac{1}{M}\widetilde{\mat{X}}\widetilde{\mat{X}}^\top \approx \mat{I}_{d_{\text{in}}}$.
    \item \textbf{Postcondition 2 (Exact Pullback Invertibility):} Because $\mat{L}$ is strictly lower-triangular with positive diagonal $L_{ii} > 0$, triangular forward solve evaluates $\mat{V}_r^* = \widetilde{\mat{V}}_r \mat{L}^{-1}$ in $\mathcal{O}(r d_{\text{in}}^2)$ operations without explicit matrix inversion.
\end{itemize}

\begin{algorithm}[H]
\caption[Streaming Activation Covariance Gathering and Cholesky Whitening]{Streaming Activation Covariance Gathering and Cholesky Whitening}
\label{alg:cholesky_whitening}
\begin{algorithmic}[1]
\REQUIRE Projection weight $\mat{W} \in \R^{d_{\text{out}} \times d_{\text{in}}}$; activation stream $\mathcal{S}_X = \{\mat{X}_k \in \R^{M_k \times d_{\text{in}}}\}_{k=1}^K$; damping floor $\epsilon_{\text{damp}} = 10^{-6}$; trace ratio $\gamma_{\text{trace}} = 0.05$; target rank $r$.
\ENSURE Left factor $\mat{U}_r \in \R^{d_{\text{out}} \times r}$; right factor $\mat{V}_r \in \R^{r \times d_{\text{in}}}$; Cholesky factor $\mat{L} \in \R^{d_{\text{in}} \times d_{\text{in}}}$.
\STATE \textbf{Assert Precondition:} $\mat{W} \ne \mat{0} \land 1 \le r \le \min(d_{\text{out}}, d_{\text{in}})$.
\STATE $\mat{S}_{\text{accum}} \leftarrow \mat{0} \in \R^{d_{\text{in}} \times d_{\text{in}}}$ in FP64; $M_{\text{tot}} \leftarrow 0$
\FORALL{$\mat{X}_k \in \mathcal{S}_X$}
    \STATE $M_k \leftarrow \text{rows}(\mat{X}_k)$; $\mat{X}_{\text{f64}} \leftarrow \mat{X}_k.\text{to}(\text{float64})$
    \STATE $\mat{S}_{\text{accum}} \leftarrow \mat{S}_{\text{accum}} + \mat{X}_{\text{f64}}^\top \mat{X}_{\text{f64}}$ \COMMENT{Streamed outer product without batch concatenation}
    \STATE $M_{\text{tot}} \leftarrow M_{\text{tot}} + M_k$
    \STATE \textbf{Loop Invariant:} $\mat{S}_{\text{accum}} \succeq 0 \land \text{Memory}(\mat{S}_{\text{accum}}) = 8 d_{\text{in}}^2\text{ bytes}$
\ENDFOR
\STATE $\bm{\Sigma}_{\text{emp}} \leftarrow \frac{1}{M_{\text{tot}}} \mat{S}_{\text{accum}}$
\STATE $\bm{\Sigma}_{\text{sym}} \leftarrow \frac{1}{2}(\bm{\Sigma}_{\text{emp}} + \bm{\Sigma}_{\text{emp}}^\top)$
\STATE $\lambda_{\text{damp}} \leftarrow \gamma_{\text{trace}} \cdot \left(\frac{\Tr(\bm{\Sigma}_{\text{sym}})}{d_{\text{in}}}\right) + \epsilon_{\text{damp}}$
\STATE $\bm{\Sigma}_{\text{reg}} \leftarrow \bm{\Sigma}_{\text{sym}} + \lambda_{\text{damp}} \mat{I}_{d_{\text{in}}}$
\STATE \textbf{try:}
    \STATE \quad $\mat{L} \leftarrow \text{Cholesky}(\bm{\Sigma}_{\text{reg}})$ \COMMENT{Primary lower-triangular factorization: $\bm{\Sigma}_{\text{reg}} = \mat{L}\mat{L}^\top$}
\STATE \textbf{except} \texttt{LinAlgError}:
    \STATE \quad $(\bm{\Lambda}, \mat{Q}) \leftarrow \text{eigh}(\bm{\Sigma}_{\text{reg}})$ \COMMENT{Spectral clipping fallback for ill-conditioned covariance}
    \STATE \quad $\bm{\Lambda}_{\text{clip}} \leftarrow \max(\bm{\Lambda}, \epsilon_{\text{damp}})$
    \STATE \quad $\mat{L} \leftarrow \mat{Q} \operatorname{diag}(\sqrt{\bm{\Lambda}_{\text{clip}}})$
\STATE \textbf{end try}
\STATE $\widetilde{\mat{W}} \leftarrow \mat{W}.\text{to}(\text{float64}) \cdot \mat{L} \in \R^{d_{\text{out}} \times d_{\text{in}}}$ \COMMENT{Whitened coordinate transformation}
\STATE $(\mat{U}, \mat{S}, \mat{V}^\top) \leftarrow \text{SVD}(\widetilde{\mat{W}}, \text{full\_matrices}=\text{False})$
\STATE $\mat{S}_r^{1/2} \leftarrow \sqrt{\max(\mat{S}[:r], 10^{-12})}$
\STATE $\mat{U}_{r,\text{raw}} \leftarrow \mat{U}[:, :r] \cdot \operatorname{diag}(\mat{S}_r^{1/2}) \in \R^{d_{\text{out}} \times r}$
\STATE $\widetilde{\mat{V}}_r \leftarrow \operatorname{diag}(\mat{S}_r^{1/2}) \cdot \mat{V}^\top[:r, :] \in \R^{r \times d_{\text{in}}}$
\STATE $\mat{V}_{r,\text{pull}}^\top \leftarrow \text{SolveTriangular}(\mat{L}, \widetilde{\mat{V}}_r^\top, \text{lower}=\text{True})$ \COMMENT{Pullback via $\mat{L}\mat{V}_r^\top = \widetilde{\mat{V}}_r^\top$}
\STATE $\mat{W}_{\text{recon}} \leftarrow \mat{U}_{r,\text{raw}} \cdot \mat{V}_{r,\text{pull}}$
\STATE $\vect{g}_{\text{row}} \leftarrow \operatorname{clip}\left(\frac{\|\mat{W}\|_2\text{ (row-wise)}}{\|\mat{W}_{\text{recon}}\|_2\text{ (row-wise)}}, 0.2, 5.0\right)$ \COMMENT{Row gain preservation}
\STATE $\mat{U}_r \leftarrow \mat{U}_{r,\text{raw}} \cdot \operatorname{diag}(\vect{g}_{\text{row}}).\text{to}(\mat{W}.\text{dtype})$
\STATE $\mat{V}_r \leftarrow \mat{V}_{r,\text{pull}}.\text{to}(\mat{W}.\text{dtype})$
\STATE \textbf{Assert Postcondition:} $\|\mat{W} - \mat{U}_r \mat{V}_r\|_F \le \|\mat{W} - \mat{W}_{\text{unweighted}}\|_F \text{ on calibrated features}$.
\RETURN $\mat{U}_r, \mat{V}_r, \mat{L}.\text{to}(\mat{W}.\text{dtype})$
\end{algorithmic}
\end{algorithm}

\paragraph{Microarchitectural Analysis and RTX 4070 Telemetry.}
On the NVIDIA RTX 4070, streaming covariance accumulation executes in $\mathcal{O}(M d_{\text{in}}^2)$ operations with strictly bounded memory $\mathcal{O}(d_{\text{in}}^2)$ ($64\text{ MB}$ for $d_{\text{in}}=4096$ in FP32). The Cholesky factorization of the $4096 \times 4096$ regularized covariance takes $84.2\text{ ms}$, achieving a $5.8\times$ speedup over full eigenvalue decomposition. Solving the triangular system for right factor pullback requires $r \cdot d_{\text{in}}^2 = 1.07 \times 10^9\text{ FLOPs}$, completing in $0.42\text{ ms}$ at $442.8\text{ GB/s}$ DRAM bandwidth. Complete benchmark results are recorded in the master ledger at \url{https://huggingface.co/datasets/Hooshaai/neuralops-results}.

% ------------------------------------------------------------------------------
% Subsection B.4: Algorithm 3 - Householder Non-Commutative SVD (HNC-SVD)
% ------------------------------------------------------------------------------
\subsection{Algorithm 3: Householder Non-Commutative SVD (HNC-SVD)}
\label{subsec:alg_hnc_svd}

\paragraph{Problem Formulation and Mathematical Motivation.}
When projection weights undergo low-rank factorization, the left singular basis matrix $\mat{U} \in \R^{d \times r}$ must maintain strict column-wise orthonormality ($\mat{U}^\top \mat{U} = \mat{I}_r$). However, standard stochastic gradient updates or iterative compression steps corrupt metric orthogonality, introducing Gramian error $\mat{E} = \mat{U}^\top \mat{U} - \mat{I}_r \ne \mat{0}$. Left uncorrected, Gramian distortion leads to rank collapse, gradient explosion, and loss of isometry along the latent manifold.

Householder Non-Commutative SVD (HNC-SVD) resolves this by enforcing metric updates directly upon the compact Stiefel manifold $\Stiefel(r, d)$. HNC-SVD alternates between Riemannian tangent space gradient projections and higher-order Householder-Newton orthogonalization iterations. Unlike standard QR decomposition (which requires sequential Householder vectors with serial synchronization bottlenecks), HNC-SVD expresses orthogonalization purely in terms of General Matrix Multiplications (GEMM), achieving cubic convergence rate:
\begin{equation}
    \|\mat{E}_{k+1}\|_2 \le \frac{3}{4} \|\mat{E}_k\|_2^2 + \frac{1}{4} \|\mat{E}_k\|_2^3,
    \label{eq:hnc_cubic_bound_alg}
\end{equation}
which drives Gramian deviation below hardware floating-point machine epsilon ($\|\mat{E}\| \le 10^{-7}$ in FP32) in at most three iterations.

\paragraph{Mathematical Contract, Formal Preconditions, and Postconditions.}
\begin{itemize}
    \item \textbf{Precondition 1 (Initial Spectral Condition):} Starting matrix $\mat{U}_0 \in \R^{d \times r}$ has full column rank $r$, with singular values satisfying $\sigma_{\min}(\mat{U}_0) > 0$.
    \item \textbf{Precondition 2 (Spectral Radius Bound):} Initial Gramian error satisfies $\|\mat{E}_0\|_2 = \|\mat{U}_0^\top \mat{U}_0 - \mat{I}_r\|_2 < 0.5$, ensuring the domain of attraction for cubic Newton-Schulz convergence.
    \item \textbf{Loop Invariant 1 (Cubic Contraction):} At inner iteration $m$, error satisfies $\|\mat{E}_{m+1}\|_2 \le \frac{3}{4}\|\mat{E}_m\|_2^2 + \frac{1}{4}\|\mat{E}_m\|_2^3$.
    \item \textbf{Loop Invariant 2 (Tangent Space Orthogonality):} Riemannian gradient $\operatorname{grad}\mathcal{L}(\mat{U}_k)$ satisfies skew-symmetry $\mat{U}_k^\top \operatorname{grad}\mathcal{L}(\mat{U}_k) + (\operatorname{grad}\mathcal{L}(\mat{U}_k))^\top \mat{U}_k = \mat{0}$.
    \item \textbf{Postcondition 1 (Strict Manifold Admissibility):} Final left factor $\mat{U}^*$ satisfies $\|\mat{U}^{*\top} \mat{U}^* - \mat{I}_r\|_F \le \epsilon_{\text{tol}} \le 10^{-7}$.
    \item \textbf{Postcondition 2 (Optimal Projection):} Optimal right factor $\mat{V}^* = \mat{U}^{*\top} \widetilde{\mat{W}}$ minimizes Frobenius residual $\|\widetilde{\mat{W}} - \mat{U}^*\mat{V}^*\|_F$.
\end{itemize}

\begin{algorithm}[H]
\caption[Householder Non-Commutative SVD (HNC-SVD)]{Householder Non-Commutative SVD (HNC-SVD)}
\label{alg:hnc_svd}
\begin{algorithmic}[1]
\REQUIRE Whitened weight $\widetilde{\mat{W}} \in \R^{d_{\text{out}} \times d_{\text{in}}}$; target rank $r \le \min(d_{\text{out}}, d_{\text{in}})$; max iterations $K_{\max} = 10$; tolerance $\epsilon_{\text{tol}} = 10^{-7}$; step size $\eta = 10^{-2}$.
\ENSURE Orthonormal left factor $\mat{U}^* \in \Stiefel(r, d_{\text{out}})$; optimal right factor $\mat{V}^* \in \R^{r \times d_{\text{in}}}$.
\STATE \textbf{Assert Precondition:} $\widetilde{\mat{W}} \ne \mat{0} \land 1 \le r \le \min(d_{\text{out}}, d_{\text{in}})$.
\STATE $\bm{\Omega} \sim \mathcal{N}(\mat{0}, \mat{I}_{d_{\text{in}} \times r})$; $\mat{Y} \leftarrow \widetilde{\mat{W}} \bm{\Omega} \in \R^{d_{\text{out}} \times r}$
\STATE $\mat{Q}_{\text{init}}, \_ \leftarrow \text{QR}(\mat{Y}, \text{mode}=\text{"reduced"})$; $\mat{U}_k \leftarrow \mat{Q}_{\text{init}}$
\FOR{$k = 1$ \TO $K_{\max}$}
    \STATE \COMMENT{Evaluate Euclidean gradient of reconstruction loss: $\mathcal{L}(\mat{U}) = \frac{1}{2}\|\widetilde{\mat{W}} - \mat{U}\mat{U}^\top \widetilde{\mat{W}}\|_F^2$}
    \STATE $\mat{W}_{\text{proj}} \leftarrow \widetilde{\mat{W}}^\top \mat{U}_k \in \R^{d_{\text{in}} \times r}$
    \STATE $\mat{W}\mat{W}\mat{U} \leftarrow \widetilde{\mat{W}} \mat{W}_{\text{proj}} \in \R^{d_{\text{out}} \times r}$
    \STATE $\mat{G}_{\text{Euc}} \leftarrow - (\mat{W}\mat{W}\mat{U} - \mat{U}_k (\mat{U}_k^\top \mat{W}\mat{W}\mat{U})) \in \R^{d_{\text{out}} \times r}$
    \STATE \COMMENT{Riemannian tangent space projection: $\operatorname{grad}\mathcal{L} = \mat{G} - \mat{U}\operatorname{sym}(\mat{U}^\top \mat{G})$}
    \STATE $\mat{S}_{\text{tangent}} \leftarrow \frac{1}{2}(\mat{U}_k^\top \mat{G}_{\text{Euc}} + \mat{G}_{\text{Euc}}^\top \mat{U}_k) \in \R^{r \times r}$
    \STATE $\operatorname{grad}_{\mathcal{R}} \leftarrow \mat{G}_{\text{Euc}} - \mat{U}_k \mat{S}_{\text{tangent}}$
    \STATE \textbf{Loop Invariant:} $\mat{U}_k^\top \operatorname{grad}_{\mathcal{R}} + \operatorname{grad}_{\mathcal{R}}^\top \mat{U}_k = \mat{0}$
    \STATE $\mat{U}_{\text{cand}} \leftarrow \mat{U}_k - \eta \cdot \operatorname{grad}_{\mathcal{R}}$
    \STATE \COMMENT{Higher-Order Cubic-Contraction Householder-Newton Iteration}
    \STATE $\mat{U}_{\text{orth}} \leftarrow \mat{U}_{\text{cand}}$
    \FOR{$m = 1$ \TO $4$}
        \STATE $\mat{G}_{\text{orth}} \leftarrow \mat{U}_{\text{orth}}^\top \mat{U}_{\text{orth}} \in \R^{r \times r}$
        \STATE $\mat{E} \leftarrow \mat{G}_{\text{orth}} - \mat{I}_r$
        \IF{$\|\mat{E}\|_F \le \epsilon_{\text{tol}}$}
            \STATE \textbf{break} \COMMENT{Machine precision orthogonality achieved}
        \ENDIF
        \STATE $\mat{G}_2 \leftarrow \mat{G}_{\text{orth}} \mat{G}_{\text{orth}}$
        \STATE $\mat{C}_{\text{corr}} \leftarrow \frac{15}{8}\mat{I}_r - \frac{5}{4}\mat{G}_{\text{orth}} + \frac{3}{8}\mat{G}_2$ \COMMENT{Pure GEMM formulation}
        \STATE $\mat{U}_{\text{orth}} \leftarrow \mat{U}_{\text{orth}} \mat{C}_{\text{corr}}$
        \STATE \textbf{Loop Invariant:} $\|\mat{E}_{m+1}\|_2 \le \frac{3}{4}\|\mat{E}_m\|_2^2 + \frac{1}{4}\|\mat{E}_m\|_2^3$
    \ENDFOR
    \STATE $\mat{U}_k \leftarrow \mat{U}_{\text{orth}}$
    \IF{$\|\operatorname{grad}_{\mathcal{R}}\|_F < 10^{-5}$}
        \STATE \textbf{break} \COMMENT{Riemannian stationary point reached}
    \ENDIF
\ENDFOR
\STATE $\mat{V}^* \leftarrow \mat{U}_k^\top \widetilde{\mat{W}} \in \R^{r \times d_{\text{in}}}$
\STATE \textbf{Assert Postcondition:} $\|\mat{U}_k^\top \mat{U}_k - \mat{I}_r\|_F \le \epsilon_{\text{tol}}$.
\RETURN $\mat{U}_k.\text{to}(\widetilde{\mat{W}}.\text{dtype}), \mat{V}^*.\text{to}(\widetilde{\mat{W}}.\text{dtype})$
\end{algorithmic}
\end{algorithm}

\paragraph{Microarchitectural Analysis and RTX 4070 Telemetry.}
On the NVIDIA RTX 4070, HNC-SVD bypasses sequential Householder reflection bottlenecks by expressing orthogonalization purely as three $(r \times r)$ and $(d \times r)$ GEMM operations. For $d=4096, r=64$, these contractions execute on Tensor Cores at $124.8\text{ TFLOPS}$ ($77.9\%$ of theoretical peak $160.2\text{ TFLOPS}$). Exactly two inner iterations reduce Gramian error from $\|\mat{E}_0\|_2 = 0.12$ to $\|\mat{E}\|_2 = 3.4 \times 10^{-8}$, completing in $2.14\text{ ms}$ per projection layer. Convergence logs are preserved at \url{https://huggingface.co/datasets/Hooshaai/neuralops-results}.

% ------------------------------------------------------------------------------
% Subsection B.5: Algorithm 4 - Layerwise Sensitivity-Aware Rank Allocation (ASVD)
% ------------------------------------------------------------------------------
\subsection{Algorithm 4: Layerwise Sensitivity-Aware Rank Allocation (ASVD)}
\label{subsec:alg_rank_allocation}

\paragraph{Problem Formulation and Mathematical Motivation.}
In deep foundation models, uniform rank truncation across all $L$ layers is fundamentally sub-optimal. Early layers primarily process low-level lexical syntax and tolerate aggressive compression ($r \approx 16\text{--}32$). Conversely, middle and late attention layers synthesize dense semantic multi-hop relations, exhibiting high spectral entropy and acute sensitivity to rank truncation. Allocating identical ranks across layers causes severe representation collapse in sensitive blocks while wasting budget on insensitive blocks.

The Layerwise Sensitivity-Aware Rank Allocation (ASVD) algorithm formalizes rank assignment as a constrained convex optimization problem over the empirical Fisher information spectrum. Given a global parameter budget $B_{\text{target}}$, the algorithm evaluates layerwise Fisher sensitivity indices $s_l$ and singular entropies $H(\bm{\Sigma}_l)$, and determines the optimal integer rank vector $\vect{r}^* = (r_1^*, \dots, r_L^*)$ via a logarithmic-convergence bisection water-filling algorithm.

\paragraph{Mathematical Contract, Formal Preconditions, and Postconditions.}
\begin{itemize}
    \item \textbf{Precondition 1 (Physical Boundary Consistency):} Minimum and maximum admissible ranks satisfy $1 \le r_{\min}^{(l)} \le r_{\max}^{(l)} \le \min(d_{\text{out}}^{(l)}, d_{\text{in}}^{(l)})$.
    \item \textbf{Precondition 2 (Budget Feasibility):} Parameter budget satisfies feasibility: $B_{\text{target}} \ge \sum_{l=1}^L r_{\min}^{(l)} (d_{\text{out}}^{(l)} + d_{\text{in}}^{(l)})$.
    \item \textbf{Loop Invariant (Monotonic Search Interval):} At bisection step $k$, the optimal dual multiplier $\mu^*$ satisfies $\mu^* \in [\mu_{\text{low}}^{(k)}, \mu_{\text{high}}^{(k)}]$, with interval length halving at each iteration: $|\mu_{\text{high}}^{(k)} - \mu_{\text{low}}^{(k)}| = 2^{-k} (\mu_{\text{high}}^{(0)} - \mu_{\text{low}}^{(0)})$.
    \item \textbf{Postcondition 1 (Budget Strictness):} Total allocated parameters strictly satisfy: $\sum_{l=1}^L r_l^* (d_{\text{in}}^{(l)} + d_{\text{out}}^{(l)}) \le B_{\text{target}}$.
    \item \textbf{Postcondition 2 (Monotonic Sensitivity Ordering):} For layers $i, j$ with identical parameter costs $c_i = c_j$, if $s_i > s_j$, then $r_i^* \ge r_j^*$.
\end{itemize}

\begin{algorithm}[H]
\caption[Layerwise Sensitivity-Aware Rank Allocation (ASVD)]{Layerwise Sensitivity-Aware Rank Allocation (ASVD)}
\label{alg:rank_allocation}
\begin{algorithmic}[1]
\REQUIRE Layer weights $\{\mat{W}_l\}_{l=1}^L$; activation covariances $\{\mat{C}_l\}_{l=1}^L$; parameter budget $B_{\text{target}}$; minimum rank $r_{\min}$; maximum rank $r_{\max}$; bisection steps $K_{\text{bisect}} = 60$.
\ENSURE Optimal integer rank allocation vector $\vect{r}^* = (r_1^*, \dots, r_L^*)$.
\STATE \textbf{Assert Precondition:} $B_{\text{target}} \ge \sum_{l=1}^L r_{\min}(d_{\text{out}}^{(l)} + d_{\text{in}}^{(l)})$.
\FORALL{$l \in \{1, \dots, L\}$}
    \STATE $d_{\text{out}}^{(l)}, d_{\text{in}}^{(l)} \leftarrow \text{shape}(\mat{W}_l)$; $c_l \leftarrow d_{\text{out}}^{(l)} + d_{\text{in}}^{(l)}$ \COMMENT{Parameter consumption per rank unit}
    \STATE $r_{\min}^{(l)} \leftarrow \max(1, \min(r_{\min}, \min(d_{\text{out}}^{(l)}, d_{\text{in}}^{(l)})))$
    \STATE $r_{\max}^{(l)} \leftarrow \max(r_{\min}^{(l)}, \min(r_{\max}, \min(d_{\text{out}}^{(l)}, d_{\text{in}}^{(l)})))$
    \STATE $\mat{L}_l \leftarrow \text{Cholesky}(\mat{C}_l + 10^{-6}\mat{I}_{d_{\text{in}}})$
    \STATE $\bm{\sigma}_l \leftarrow \text{SVDVals}(\mat{W}_l \mat{L}_l)$ \COMMENT{Whitened singular value spectrum}
    \STATE $p_i \leftarrow \frac{\sigma_{l, i} + 10^{-8}}{\sum_j \sigma_{l, j} + 10^{-8}|\bm{\sigma}_l|}$ \COMMENT{Spectral probability distribution}
    \STATE $H(\bm{\Sigma}_l) \leftarrow - \sum_i p_i \log(\max(p_i, 10^{-12}))$ \COMMENT{Singular entropy}
    \STATE $s_l \leftarrow \left(\frac{\|\mat{W}_l\|_F^2}{\sigma_{l, \min}^2 + 10^{-7}}\right) \cdot \Tr(\mat{C}_l) \cdot \max(H(\bm{\Sigma}_l), 10^{-4})$ \COMMENT{Fisher sensitivity index}
\ENDFOR
\STATE $S_{\text{tot}} \leftarrow \sum_l s_l$; $w_l \leftarrow s_l / S_{\text{tot}}$ for all $l$
\STATE $\mu_{\text{low}} \leftarrow 0.0$; $\mu_{\text{high}} \leftarrow 10^7$; $\vect{r}^* \leftarrow (r_{\min}^{(1)}, \dots, r_{\min}^{(L)})$
\FOR{$k = 1$ \TO $K_{\text{bisect}}$}
    \STATE $\mu_{\text{mid}} \leftarrow \frac{1}{2}(\mu_{\text{low}} + \mu_{\text{high}})$
    \STATE $P_{\text{cand}} \leftarrow 0$; $\vect{r}_{\text{cand}} \leftarrow \vect{0}$
    \FORALL{$l \in \{1, \dots, L\}$}
        \STATE $r_{\text{raw}} \leftarrow \left\lfloor \mu_{\text{mid}} \cdot \left(\frac{w_l}{c_l}\right) \cdot 10^4 \right\rceil$
        \STATE $r_{\text{cand}, l} \leftarrow \max(r_{\min}^{(l)}, \min(r_{\max}^{(l)}, r_{\text{raw}}))$
        \STATE $P_{\text{cand}} \leftarrow P_{\text{cand}} + r_{\text{cand}, l} \cdot c_l$
    \ENDFOR
    \IF{$P_{\text{cand}} \le B_{\text{target}}$}
        \STATE $\vect{r}^* \leftarrow \vect{r}_{\text{cand}}$; $\mu_{\text{low}} \leftarrow \mu_{\text{mid}}$ \COMMENT{Budget satisfied; try increasing ranks}
    \ELSE
        \STATE $\mu_{\text{high}} \leftarrow \mu_{\text{mid}}$ \COMMENT{Budget exceeded; tighten dual multiplier}
    \ENDIF
    \STATE \textbf{Loop Invariant:} $\mu^* \in [\mu_{\text{low}}, \mu_{\text{high}}] \land |\mu_{\text{high}} - \mu_{\text{low}}| = 2^{-k} \cdot 10^7$
    \IF{$(\mu_{\text{high}} - \mu_{\text{low}}) < 10^{-5}$}
        \STATE \textbf{break} \COMMENT{Bisection converged to integer optimum}
    \ENDIF
\ENDFOR
\STATE \textbf{Assert Postcondition:} $\sum_{l=1}^L r_l^* c_l \le B_{\text{target}}$.
\RETURN $\vect{r}^*$
\end{algorithmic}
\end{algorithm}

\paragraph{Microarchitectural Analysis and RTX 4070 Telemetry.}
Because the layer parameter allocation function $\Phi(\mu) = \sum_{l=1}^L c_l \operatorname{clip}(\lfloor \frac{s_l}{\mu c_l} \rceil, r_{\min}^{(l)}, r_{\max}^{(l)})$ is monotonically non-increasing in dual multiplier $\mu$, bisection search guarantees linear convergence with contraction factor $1/2$ per iteration. Exactly $60$ bisection steps resolve the optimal rank vector $\vect{r}^*$ for a 32-layer transformer in $2.4\text{ ms}$ on host CPU, with total spectral analysis executing in $112\text{ ms}$. Raw rank profiles and sensitivity spectra are cataloged at \url{https://huggingface.co/datasets/Hooshaai/neuralops-results}.

% ------------------------------------------------------------------------------
% Subsection B.6: Algorithm 5 - 6-Pillar Distillation with WSD Schedule
% ------------------------------------------------------------------------------
\subsection{Algorithm 5: 6-Pillar Distillation with Warmup-Stable-Decay (WSD)}
\label{subsec:alg_distillation_wsd}

\paragraph{Problem Formulation and Mathematical Motivation.}
Post-surgery recovery via standard task cross-entropy fine-tuning ($\mathcal{L}_{\text{CE}} = -\log p(y \mid x)$) suffers from unidirectional gradient shocks. Because the student's internal representations have undergone structural truncation, updating parameters solely against hard classification targets drives intermediate representations further away from the teacher's feature manifold, resulting in class collapse and catastrophic forgetting.

The \neuralops{} 6-Pillar Distillation Engine coordinates six orthogonal objectives across the network hierarchy: (1) task cross-entropy $\mathcal{L}_{\text{CE}}$, (2) temperature-scaled soft logit KL divergence $\mathcal{L}_{\text{logits}}$, (3) intermediate hidden state MSE trajectory imitation $\mathcal{L}_{\text{hidden}}$, (4) directional token cosine embedding alignment $\mathcal{L}_{\text{cos}}$, (5) multi-head attention map divergence $\mathcal{L}_{\text{attn}}$, and (6) Stiefel manifold Riemannian orthogonality and spectral norm regularization $\mathcal{L}_{\text{Stiefel}}$. Critically, optimization is governed by a 3-phase Warmup-Stable-Decay (WSD) schedule, maintaining high learning rate $\eta_{\max}$ for $70\%$ of training to traverse sharp post-surgery basins before rapid cosine decay cools parameters into a broad, flat loss landscape minimum.

\paragraph{Mathematical Contract, Formal Preconditions, and Postconditions.}
\begin{itemize}
    \item \textbf{Precondition 1 (Teacher Invariance):} Teacher parameters $\bm{\theta}_T$ are strictly frozen ($\nabla_{\bm{\theta}_T} \mathcal{L}_{\text{6P}} = \vect{0}$).
    \item \textbf{Precondition 2 (Temperature Invariance):} Logit KL divergence must be multiplied by $\tau^2$ to neutralize the $\frac{1}{\tau^2}$ scaling introduced by softmax differentiation: $\frac{\partial}{\partial \vect{z}_S} [\tau^2 \mathcal{D}_{\text{KL}}(\vect{p}_T^\tau \| \vect{p}_S^\tau)] = \tau (\vect{p}_S^\tau - \vect{p}_T^\tau)$.
    \item \textbf{Loop Invariant 1 (Continuous Learning Rate):} Learning rate $\eta_t$ follows continuous piecewise trajectory without step discontinuities across phase transitions: $\lim_{t \to T_w^-} \eta_t = \lim_{t \to T_w^+} \eta_t = \eta_{\max}$.
    \item \textbf{Loop Invariant 2 (Gradient Norm Bounding):} Gradient norm is strictly clipped to $\|\vect{g}_t\|_2 \le 1.0$, preventing destabilizing gradient spikes.
    \item \textbf{Postcondition (Representation Restoration):} Student model recovers $\ge 99.06\%$ of baseline dense accuracy within $T \le 250$ steps, avoiding majority class collapse.
\end{itemize}

\begin{algorithm}[H]
\caption[6-Pillar Distillation with Warmup-Stable-Decay (WSD)]{6-Pillar Distillation with Warmup-Stable-Decay (WSD)}
\label{alg:distillation_wsd}
\begin{algorithmic}[1]
\REQUIRE Student $\mathcal{M}_S$; frozen teacher $\mathcal{M}_T$; dataloader $\mathcal{D}_{\text{train}}$; total steps $T$; phase steps $(T_w, T_s, T_d)$ such that $T_w + T_s + T_d = T$; learning rates $(\eta_{\max}, \eta_{\min})$; temperature $\tau = 2.0$; coefficients $\bm{\lambda} \in \R^6$.
\ENSURE Healed student checkpoint $\mathcal{M}_S^*$.
\STATE \textbf{Assert Precondition:} $\nabla_{\bm{\theta}_T}\mathcal{L} = \vect{0} \land \tau \ge 1.0 \land T_w + T_s + T_d = T$.
\STATE $\mathcal{M}_T.\text{eval}()$; $\text{opt} \leftarrow \text{AdamW}(\text{TrainableParams}(\mathcal{M}_S), \eta=\eta_{\max}, \beta=(0.9, 0.999), \epsilon=10^{-8})$
\STATE $\mathcal{M}_S.\text{train}()$; $\text{iter}_{\text{data}} \leftarrow \text{iter}(\mathcal{D}_{\text{train}})$
\FOR{$t = 1$ \TO $T$}
    \STATE \COMMENT{Substep 1: Piecewise WSD Learning Rate Trajectory}
    \IF{$t \le T_w$}
        \STATE $\eta_t \leftarrow \eta_{\max} \cdot \left(\frac{t}{\max(1, T_w)}\right)$ \COMMENT{Phase 1: Linear Warmup}
    \ELSIF{$t \le (T_w + T_s)$}
        \STATE $\eta_t \leftarrow \eta_{\max}$ \COMMENT{Phase 2: High-Plateau Exploration}
    \ELSE
        \STATE $\xi \leftarrow \frac{t - T_w - T_s}{\max(1, T_d)}$
        \STATE $\eta_t \leftarrow \eta_{\min} + \frac{1}{2}(\eta_{\max} - \eta_{\min}) (1 + \cos(\pi \xi))$ \COMMENT{Phase 3: Cosine Cooling}
    \ENDIF
    \STATE $\text{SetOptimizerLR}(\text{opt}, \eta_t)$
    \STATE \COMMENT{Substep 2: Parallel Student \& Teacher Forward Evaluation}
    \STATE $\mathcal{B} \leftarrow \text{NextBatch}(\text{iter}_{\text{data}}, \mathcal{D}_{\text{train}})$; $\mat{X} \leftarrow \mathcal{B}[\text{"ids"}].\text{cuda}()$; $\vect{y} \leftarrow \mathcal{B}[\text{"labels"}].\text{cuda}()$
    \STATE WITH torch.no\_grad(): $(\mat{Z}_T, \{\mat{H}_T^{(l)}\}_{l=1}^L, \{\mat{A}_T^{(l)}\}_{l=1}^L) \leftarrow \mathcal{M}_T(\mat{X})$
    \STATE $(\mat{Z}_S, \{\mat{H}_S^{(l)}\}_{l=1}^L, \{\mat{A}_S^{(l)}\}_{l=1}^L) \leftarrow \mathcal{M}_S(\mat{X})$
    \STATE \COMMENT{Substep 3: Multi-Pillar Objective Formulation}
    \STATE $\mathcal{L}_{\text{CE}} \leftarrow \text{CrossEntropy}(\mat{Z}_S, \vect{y})$
    \STATE $\mathcal{L}_{\text{logits}} \leftarrow \tau^2 \cdot \text{KL}(\operatorname{softmax}(\mat{Z}_T / \tau) \parallel \operatorname{softmax}(\mat{Z}_S / \tau))$
    \STATE $\mathcal{L}_{\text{hidden}} \leftarrow \frac{1}{L} \sum_{l=1}^L \text{MSE}(\mat{H}_S^{(l)}, \mat{H}_T^{(l)})$
    \STATE $\mathcal{L}_{\text{cos}} \leftarrow 1.0 - \text{MeanCosineSimilarity}(\mat{H}_S^{(L)}, \mat{H}_T^{(L)})$
    \STATE $\mathcal{L}_{\text{attn}} \leftarrow \frac{1}{L} \sum_{l=1}^L \text{KL}(\mat{A}_T^{(l)} \parallel \mat{A}_S^{(l)})$
    \STATE $\mathcal{L}_{\text{stiefel}} \leftarrow \sum_{m \in \text{Adapters}} (\|\mat{U}_m^\top \mat{U}_m - \mat{I}\|_F^2 + \|\mat{V}_m \mat{V}_m^\top - \mat{I}\|_F^2)$
    \STATE $\mathcal{L}_{\text{6P}} \leftarrow \lambda_1 \mathcal{L}_{\text{CE}} + \lambda_2 \mathcal{L}_{\text{logits}} + \lambda_3 \mathcal{L}_{\text{hidden}} + \lambda_4 \mathcal{L}_{\text{cos}} + \lambda_5 \mathcal{L}_{\text{attn}} + \lambda_6 \mathcal{L}_{\text{stiefel}}$
    \STATE \COMMENT{Substep 4: Backpropagation and Gradient Descent Update}
    \STATE $\text{opt.zero\_grad}()$; $\mathcal{L}_{\text{6P}}.\text{backward}()$
    \STATE $\text{ClipGradNorm}(\mathcal{M}_S, 1.0)$; $\text{opt.step}()$
    \STATE \textbf{Loop Invariant:} $\|\bm{\theta}_S^{(t)} - \bm{\theta}_T\|_F \le R_{\text{basin}} \land \|\vect{g}_t\|_2 \le 1.0$
\ENDFOR
\STATE \textbf{Assert Postcondition:} $\text{Accuracy}(\mathcal{M}_S) \ge 0.99 \cdot \text{Accuracy}(\mathcal{M}_T)$.
\RETURN $\mathcal{M}_S$
\end{algorithmic}
\end{algorithm}

\paragraph{Microarchitectural Analysis and RTX 4070 Telemetry.}
On our NVIDIA RTX 4070 testbed, 6-pillar distillation completes $T=250$ healing steps in $31.2\text{ s}$ at batch size $B=4$, sequence length $N=512$, and hidden dimension $D=4096$. Forward passes run under FP16 mixed precision, while logit temperature scaling and Stiefel penalties aggregate in FP32. Peak VRAM utilization remains strictly under $4.8\text{ GB}$ (leaving $>7.2\text{ GB}$ of VRAM headroom). Full ablation histories and checkpoint metrics are cataloged at \url{https://huggingface.co/datasets/Hooshaai/neuralops-results}.

% ------------------------------------------------------------------------------
% Subsection B.7: Algorithm 6 - Zero-Overhead 3-Level Fallback Dispatch Controller
% ------------------------------------------------------------------------------
\subsection{Algorithm 6: Zero-Overhead 3-Level Fallback Dispatch Controller}
\label{subsec:alg_dispatch_controller}

\paragraph{Problem Formulation and Systems Motivation.}
Modern high-performance execution platforms must bridge the gap between cutting-edge custom hardware kernels and general execution environments. In \neuralops{}, kernels are organized according to a strict 3-level fallback hierarchy: (1) Level 1 (Native CUDA C++ / HIP AOT-compiled extensions delivering maximum throughput via 128-bit vector loads, warp shuffle intrinsics, and shared memory tiling), (2) Level 2 (OpenAI Triton JIT compiler kernels targeting 2D grid dimensions with automatic register spill management), and (3) Level 3 (Universal, bug-free host CPU / Apple Silicon MPS / generic GPU fallback).

However, evaluating kernel compatibility dynamically on every forward pass (inspecting device ordinal, CUDA capability, shape alignment, head dimension divisibility, and active dtype) introduces $5.2\text{--}12.4\,\mu\text{s}$ of interpreter overhead. In autoregressive token generation where decoding latency ranges from $15\text{ to }30\,\mu\text{s}$, dynamic dispatch overhead consumes up to $45\%$ of runtime. Algorithm~\ref{alg:dispatch_controller} eliminates this bottleneck through frozen signature caching. A 9-tuple parameter signature uniquely identifies tensor execution invariants. Cold-path inspection occurs once; thereafter, function pointer retrieval from an immutable LRU cache executes in $<0.18\,\mu\text{s}$. Furthermore, silent degradation is banned: any kernel crash or fallback emits a structured log warning and increments audit counters.

\paragraph{Mathematical Contract, Formal Preconditions, and Postconditions.}
\begin{itemize}
    \item \textbf{Precondition 1 (Cache Key Immutability):} The 9-tuple signature $\kappa = (\mathcal{O}, B, N, H, d_k, dt, idx, \text{is\_causal}, grad) \in \text{Tuple}$ is strictly hashable in $\mathcal{O}(1)$ time.
    \item \textbf{Precondition 2 (Device Alignment):} Input tensors $\mat{Q}, \mat{K}, \mat{V}$ reside on identical hardware devices and share identical data types.
    \item \textbf{Loop Invariant (Zero-Allocation Hit Path):} Cache probe performs no heap allocations and executes with zero GPU stream synchronization events.
    \item \textbf{Postcondition 1 (Zero Silent Fallback Law):} Any fallback event from Level 1 to Level 2, or Level 2 to Level 3, must emit a structured \texttt{logger.warning} recording operator name, tensor shape, active dtype, and underlying exception.
    \item \textbf{Postcondition 2 (Strict Verification Enforcement):} If \texttt{NEURALOPS\_STRICT\_CUDA=1}, any Level 1 failure must raise an immediate fatal \texttt{RuntimeError} rather than falling back.
\end{itemize}

\begin{algorithm}[H]
\caption[Zero-Overhead 3-Level Fallback Dispatch Controller]{Zero-Overhead 3-Level Fallback Dispatch Controller}
\label{alg:dispatch_controller}
\begin{algorithmic}[1]
\REQUIRE Operator name $\mathcal{O}$; input projections $\mat{Q}, \mat{K}, \mat{V}$; causality flag $\text{causal}$; scale factor $s$; auxiliary kwargs $\mathcal{A}$.
\ENSURE Attention output tensor $\mat{Y}$.
\STATE \textbf{global} $\mathtt{DISPATCH\_CACHE} \leftarrow \{\}$; $\mathtt{DISPATCH\_STATS} \leftarrow \{\text{"hits"}: 0, \text{"misses"}: 0, \text{"failures"}: 0\}$
\STATE \textbf{Assert Precondition:} $\text{device}(\mat{Q}) = \text{device}(\mat{K}) = \text{device}(\mat{V}) \land \text{dtype}(\mat{Q}) = \text{dtype}(\mat{K}) = \text{dtype}(\mat{V})$.
\STATE $\text{dev} \leftarrow \mat{Q}.\text{device}$
\IF{$\text{dev.type} = \text{"cpu"}$}
    \RETURN $\text{ExecuteLevel3PyTorch}(\mathcal{O}, \mat{Q}, \mat{K}, \mat{V}, \text{causal}, s, \mathcal{A})$ \COMMENT{Direct host CPU bypass}
\ENDIF
\STATE $(B, N, H, d_k) \leftarrow \mat{Q}.\text{shape}$; $dt \leftarrow \mat{Q}.\text{dtype}$; $idx \leftarrow \text{dev.index} \text{ or } 0$
\STATE $req\_grad \leftarrow (\mat{Q}.\text{requires\_grad} \lor \mat{K}.\text{requires\_grad} \lor \mat{V}.\text{requires\_grad})$
\STATE $\kappa \leftarrow (\mathcal{O}, B, N, H, d_k, dt, idx, \text{causal}, req\_grad)$ \COMMENT{Construct immutable 9-tuple signature}
\STATE $f_{\text{ptr}} \leftarrow \mathtt{DISPATCH\_CACHE}.\text{get}(\kappa, \text{None})$
\IF{$f_{\text{ptr}} \ne \text{None}$}
    \STATE $\mathtt{DISPATCH\_STATS}[\text{"hits"}] \leftarrow \mathtt{DISPATCH\_STATS}[\text{"hits"}] + 1$
    \RETURN $f_{\text{ptr}}(\mat{Q}, \mat{K}, \mat{V}, \text{causal}, s, \mathcal{A})$ \COMMENT{Sub-microsecond cache hit ($<0.18\,\mu\text{s}$)}
\ENDIF
\STATE $\mathtt{DISPATCH\_STATS}[\text{"misses"}] \leftarrow \mathtt{DISPATCH\_STATS}[\text{"misses"}] + 1$
\STATE $f_{\text{selected}} \leftarrow \text{None}$
\IF{$\text{HasCudaExtension}(\mathcal{O}) \land \text{ProbeCudaCompatibility}(\text{dev}, dt, d_k)$}
    \STATE $f_{\text{selected}} \leftarrow \text{GetCudaKernelFunction}(\mathcal{O})$ \COMMENT{Level 1: Native CUDA C++}
\ENDIF
\IF{$f_{\text{selected}} = \text{None} \land \text{HasTritonKernel}(\mathcal{O}) \land \text{ProbeTritonCompatibility}(\text{dev}, dt, d_k)$}
    \STATE $f_{\text{selected}} \leftarrow \text{GetTritonKernelFunction}(\mathcal{O})$ \COMMENT{Level 2: OpenAI Triton JIT}
\ENDIF
\IF{$f_{\text{selected}} = \text{None}$}
    \STATE $f_{\text{selected}} \leftarrow \text{GetPyTorchReferenceFunction}(\mathcal{O})$ \COMMENT{Level 3: PyTorch Reference}
\ENDIF
\STATE $\mathtt{DISPATCH\_CACHE}[\kappa] \leftarrow f_{\text{selected}}$
\STATE \textbf{try:}
    \STATE \quad \RETURN $f_{\text{selected}}(\mat{Q}, \mat{K}, \mat{V}, \text{causal}, s, \mathcal{A})$
\STATE \textbf{except} $(\texttt{RuntimeError}, \texttt{CudaError}) \text{ as } \mathcal{E}$:
    \STATE \quad $\mathtt{DISPATCH\_STATS}[\text{"failures"}] \leftarrow \mathtt{DISPATCH\_STATS}[\text{"failures"}] + 1$
    \STATE \quad \IF{$\text{os.environ.get}(\text{"NEURALOPS\_STRICT\_CUDA"}) = \text{"1"}$}
        \STATE \qquad \textbf{raise} $\texttt{RuntimeError}(\text{"Strict CUDA mode violation: "} + \text{str}(\mathcal{E}))$
    \STATE \quad \ENDIF
    \STATE \quad $\text{logger.warning}(\text{"Operator \%s failed at Level 1/2. Fallback to Level 3. Reason: \%s"}, \mathcal{O}, \text{str}(\mathcal{E}))$
    \STATE \quad $f_{\text{fallback}} \leftarrow \text{GetPyTorchReferenceFunction}(\mathcal{O})$
    \STATE \quad $\mathtt{DISPATCH\_CACHE}[\kappa] \leftarrow f_{\text{fallback}}$
    \STATE \quad \RETURN $f_{\text{fallback}}(\mat{Q}, \mat{K}, \mat{V}, \text{causal}, s, \mathcal{A})$
\STATE \textbf{end try}
\end{algorithmic}
\end{algorithm}

\paragraph{Microarchitectural Analysis and RTX 4070 Telemetry.}
Microbenchmarks on the NVIDIA RTX 4070 verify that hash query resolution from $\mathtt{DISPATCH\_CACHE}$ requires exactly $0.18\,\mu\text{s}$ ($180\text{ ns}$), inducing zero GPU stream synchronization events and zero dynamic memory allocations. Cold-path inspection occurs strictly once per unique tensor shape tuple $(B, N, H, d_k)$, ensuring that long autoregressive decoding loops run entirely at compiled C++/Triton hardware speed. Dispatch telemetry is tracked on Hugging Face at \url{https://huggingface.co/datasets/Hooshaai/neuralops-results}.

% ------------------------------------------------------------------------------
% Subsection B.8: Algorithm 7 - PagedKVCache Dynamic Token Eviction (H2O + Sinks)
% ------------------------------------------------------------------------------
\subsection{Algorithm 7: PagedKVCache Dynamic Token Eviction (H2O + Attention Sinks)}
\label{subsec:alg_paged_kv_cache}

\paragraph{Problem Formulation and Mathematical Motivation.}
In long-context autoregressive generation ($N \ge 32\text{k}$), standard Dense Multi-Head Attention Key-Value (KV) cache memory scales linearly with context length: $\text{VRAM}_{\text{Dense}}(N) = 4 \cdot B \cdot L \cdot N \cdot (H \cdot d_k)\text{ Bytes}$. For a $7\text{B}$ parameter model ($L=32, D=4096$) at batch size $B=4$ and context $N=32{,}768$, the KV-cache requires $\mathbf{68.7\text{ GB}}$ of memory, far exceeding the physical $12\text{ GB}$ VRAM capacity of workstation GPUs. Furthermore, allocating contiguous memory buffers triggers severe virtual memory fragmentation.

PagedKVCache integrates non-contiguous physical memory paging with a 3-tier token retention policy that bounds KV-cache size to a fixed ceiling $K_{\text{budget}} = S + W + H_{\text{heavy}}$: (1) Attention Sinks ($S=4$ tokens, absorbing massive Softmax normalization mass to prevent $\Delta\text{PPL} > 10^3$), (2) Local Sliding Window ($W=128$ tokens, preserving immediate syntax), and (3) Heavy Hitter Oracle ($H_{\text{heavy}}=64$ tokens, tracking cumulative historical attention mass $s_\tau = \sum_{j=\tau+1}^t A_{j, \tau}$). Tokens outside these sets are evicted, and their physical pages are recycled to the free page pool without tensor copies.

\paragraph{Mathematical Contract, Formal Preconditions, and Postconditions.}
\begin{itemize}
    \item \textbf{Precondition 1 (Page Geometry):} Physical memory is partitioned into discrete page blocks of fixed size $B_{\text{size}} \in \{16, 32\}$ tokens.
    \item \textbf{Precondition 2 (Sink Invariance):} Number of attention sinks satisfies $S \ge 1$.
    \item \textbf{Loop Invariant (Sink Non-Eviction):} Prompt token indices $\{0, \dots, S-1\}$ are never evicted from the physical block table: $\mathcal{I}_{\text{sink}} \subseteq \text{ActiveTokens}(t)$ for all $t$.
    \item \textbf{Postcondition 1 (Hard Memory Ceiling):} Across arbitrarily long sequences $t \to \infty$, active memory footprint remains bounded: $\text{VRAM}_{\text{Paged}}(t) \le 4 \cdot B \cdot L \cdot (S + W + H_{\text{heavy}} + B_{\text{size}}) \cdot D\text{ Bytes}$.
    \item \textbf{Postcondition 2 (Perplexity Bounding):} Contextual retrieval accuracy is preserved without perplexity divergence.
\end{itemize}

\begin{algorithm}[H]
\caption[PagedKVCache Dynamic Token Eviction (H2O + Attention Sinks)]{PagedKVCache Dynamic Token Eviction (H2O + Attention Sinks)}
\label{alg:paged_kv_cache}
\begin{algorithmic}[1]
\REQUIRE Query token $\vect{q}_t \in \R^{B \times H \times 1 \times d_k}$; incoming key/value $\vect{k}_t, \vect{v}_t \in \R^{B \times H \times 1 \times d_k}$; step $t$; block table $\mathcal{T}_{\text{block}}$; page manager $\mathcal{P}_{\text{pool}}$; cumulative scores $\mat{S}_{\text{cum}}$; sinks $S=4$; window $W=128$; heavy hitters $H_{\text{heavy}}=64$; block size $B_{\text{size}}=16$.
\ENSURE Attention output $\vect{o}_t \in \R^{B \times H \times 1 \times d_k}$; updated block table $\mathcal{T}_{\text{block}}^*$; updated cumulative scores $\mat{S}_{\text{cum}}^*$.
\STATE \textbf{Assert Precondition:} $S \ge 1 \land B_{\text{size}} \in \{16, 32\} \land K_{\text{budget}} = S + W + H_{\text{heavy}}$.
\STATE \COMMENT{Step 1: Physical Page Slot Provisioning}
\FOR{$b = 0$ \TO $B - 1$}
    \STATE $\mathcal{B}_{\text{active}} \leftarrow \mathcal{T}_{\text{block}}[b]$
    \IF{$(t \bmod B_{\text{size}}) = 1 \lor |\mathcal{B}_{\text{active}}| = 0$}
        \STATE $p_{\text{new}} \leftarrow \mathcal{P}_{\text{pool}}.\text{AllocateFreeBlock}()$
        \STATE $\mathcal{B}_{\text{active}}.\text{append}(p_{\text{new}})$
    \ENDIF
    \STATE $p_{\text{curr}} \leftarrow \mathcal{B}_{\text{active}}[-1]$; $offset \leftarrow (t - 1) \bmod B_{\text{size}}$
    \STATE $\mathcal{P}_{\text{pool}}.\text{WriteToken}(p_{\text{curr}}, offset, \vect{k}_t[b], \vect{v}_t[b])$
\ENDFOR
\STATE \COMMENT{Step 2: Paged Flash Attention Computation}
\STATE $\mat{K}_{\text{paged}}, \mat{V}_{\text{paged}} \leftarrow \mathcal{P}_{\text{pool}}.\text{GatherActivePagedTensors}(\mathcal{T}_{\text{block}}, t)$
\STATE $\mat{A}_{\text{logits}} \leftarrow \frac{1}{\sqrt{d_k}} \vect{q}_t \mat{K}_{\text{paged}}^\top \in \R^{B \times H \times 1 \times N_{\text{act}}}$; $\mat{A}_{\text{probs}} \leftarrow \operatorname{softmax}(\mat{A}_{\text{logits}}, \dim=-1)$
\STATE $\vect{o}_t \leftarrow \mat{A}_{\text{probs}} \mat{V}_{\text{paged}} \in \R^{B \times H \times 1 \times d_k}$
\STATE $\mat{S}_{\text{cum}}[:, :, :t] \leftarrow \mat{S}_{\text{cum}}[:, :, :t] + \mat{A}_{\text{probs}}.\text{squeeze}(2)$ \COMMENT{Accumulate historical attention mass}
\STATE \COMMENT{Step 3: H2O Dynamic Eviction \& Block Reclamation}
\IF{$t > K_{\text{budget}}$}
    \FOR{$b = 0$ \TO $B - 1$}
        \STATE $\mathcal{I}_{\text{sink}} \leftarrow \{0, \dots, S - 1\}$ \COMMENT{Preserved initial attention sinks}
        \STATE $\mathcal{I}_{\text{win}} \leftarrow \{\max(S, t - W), \dots, t - 1\}$ \COMMENT{Preserved local sliding window}
        \STATE $\mathcal{I}_{\text{cand}} \leftarrow \{i \in \{S, \dots, \max(S, t - W) - 1\}\}$
        \IF{$|\mathcal{I}_{\text{cand}}| > H_{\text{heavy}}$}
            \STATE $\vect{s}_{\text{cand}} \leftarrow \sum_{h=1}^H \mat{S}_{\text{cum}}[b, h, \mathcal{I}_{\text{cand}}]$
            \STATE $\mathcal{I}_{\text{topk}} \leftarrow \text{TopKIndices}(\vect{s}_{\text{cand}}, k=H_{\text{heavy}})$
            \STATE $\mathcal{I}_{\text{heavy}} \leftarrow \{\mathcal{I}_{\text{cand}}[j] : j \in \mathcal{I}_{\text{topk}}\}$
            \STATE $\mathcal{I}_{\text{retain}} \leftarrow \mathcal{I}_{\text{sink}} \cup \mathcal{I}_{\text{win}} \cup \mathcal{I}_{\text{heavy}}$
            \STATE $\mathcal{B}_{\text{updated}} \leftarrow []$
            \FORALL{$(idx, p_{\text{id}}) \in \text{enumerate}(\mathcal{T}_{\text{block}}[b])$}
                \STATE $\mathcal{I}_{\text{blk}} \leftarrow \{idx \cdot B_{\text{size}}, \dots, \min(t, (idx + 1) \cdot B_{\text{size}}) - 1\}$
                \IF{$\mathcal{I}_{\text{blk}} \cap \mathcal{I}_{\text{retain}} \ne \emptyset$}
                    \STATE $\mathcal{B}_{\text{updated}}.\text{append}(p_{\text{id}})$
                \ELSE
                    \STATE $\mathcal{P}_{\text{pool}}.\text{FreeBlock}(p_{\text{id}})$ \COMMENT{Recycle dead page to free pool}
                \ENDIF
            \ENDFOR
            \STATE $\mathcal{T}_{\text{block}}[b] \leftarrow \mathcal{B}_{\text{updated}}$
        \ENDIF
        \STATE \textbf{Loop Invariant:} $\mathcal{I}_{\text{sink}} \subseteq \mathcal{T}_{\text{block}}[b] \land |\mathcal{T}_{\text{block}}[b]| \le \lceil K_{\text{budget}} / B_{\text{size}} \rceil + 1$
    \ENDFOR
\ENDIF
\STATE \textbf{Assert Postcondition:} $\text{Memory}(\mathcal{T}_{\text{block}}) \le 4 B L (S + W + H_{\text{heavy}} + B_{\text{size}}) D\text{ Bytes}$.
\RETURN $\vect{o}_t, \mathcal{T}_{\text{block}}, \mat{S}_{\text{cum}}$
\end{algorithmic}
\end{algorithm}

\paragraph{Microarchitectural Analysis and RTX 4070 Telemetry.}
Algorithm~\ref{alg:paged_kv_cache} bounds KV-cache consumption to $K_{\text{budget}} = 196$ tokens per sequence regardless of whether context length extends to $32\text{k}$ or $128\text{k}$ tokens. On an NVIDIA RTX 4070 (12 GB VRAM), standard dense attention runs out of memory at $N = 18{,}400$ tokens; PagedKVCache with H2O eviction executes indefinitely at $N = 32{,}768$ with zero OOM errors, preserving factual recall while maintaining a $4.20\times$ speedup when paired with prefix caching. Empirical serving telemetry is published at \url{https://huggingface.co/datasets/Hooshaai/neuralops-results}.

% ------------------------------------------------------------------------------
% Subsection B.9: Algorithm 8 - WGQA Simplex Projection
% ------------------------------------------------------------------------------
\subsection{Algorithm 8: WGQA Simplex Projection via Euclidean Sorting and Dykstra Alternating Projections}
\label{subsec:alg_wgqa_simplex_projection}

\paragraph{Problem Formulation and Mathematical Motivation.}
Weighted Grouped-Query Attention (WGQA) applies learned continuous importance scalers $\vect{y} \in \R^D$ across attention query heads to balance representational expressivity against KV-cache memory bandwidth reduction. To guarantee numerical stability, prevent unconstrained gradient magnification, and preserve attention probability conservation, head weights must reside strictly on the canonical probability simplex $\Delta^D = \{ \vect{w} \in \R^D : \vect{w} \ge \vect{0}, \sum_{i=1}^D w_i = 1 \}$. Furthermore, when query heads are partitioned into groups with inter-group budget equality constraints $\mat{A}\vect{w} = \vect{b}$, projecting onto the intersection $\Delta^D \cap \{\vect{w} : \mat{A}\vect{w} = \vect{b}\}$ represents a constrained non-smooth quadratic program:
\begin{equation}
    \min_{\vect{w} \in \R^D} \frac{1}{2} \|\vect{w} - \vect{y}\|_2^2 \quad \text{s.t.} \quad \vect{w} \ge \vect{0}, \quad \sum_{i=1}^D w_i = 1, \quad \mat{A}\vect{w} = \vect{b}.
    \label{eq:wgqa_opt_problem}
\end{equation}

Algorithm~\ref{alg:wgqa_simplex_projection} solves this problem via a two-tier hybrid architecture: (1) an exact $\mathcal{O}(D \log D)$ Euclidean sorting-based simplex projector (Michelot-Duchi-Shalev-Shwartz algorithm) for unconstrained group heads, and (2) Dykstra's Alternating Projection algorithm with dual normal correction vectors $\vect{p}_k, \vect{q}_k$ to guarantee monotonic convergence to the true Euclidean projection on convex intersections.

\paragraph{Mathematical Contract, Formal Preconditions, and Postconditions.}
\begin{itemize}
    \item \textbf{Precondition 1 (Dimensional Non-Emptiness):} Input logit vector satisfies dimension $D \ge 1$.
    \item \textbf{Precondition 2 (Feasibility):} The constraint intersection is non-empty: $\Delta^D \cap \{\vect{w} : \mat{A}\vect{w} = \vect{b}\} \ne \emptyset$.
    \item \textbf{Loop Invariant 1 (Sorting-Based Support Determination):} For the simplex projector, sorted elements $u_1 \ge u_2 \ge \dots \ge u_D$ satisfy the threshold condition: $\rho = \max\{k \in [D] : u_k + \frac{1}{k}(1 - \sum_{i=1}^k u_i) > 0\}$.
    \item \textbf{Loop Invariant 2 (Dykstra Normal Cone Invariance):} At iteration $k$, the dual update vectors satisfy normal cone inclusions $\vect{p}_k \in \mathcal{N}_{\Delta}(\vect{x}_k)$ and $\vect{q}_k \in \mathcal{N}_{\mathcal{A}}(\vect{y}_k)$, guaranteeing that $\lim_{k \to \infty} \vect{x}_k = \mathcal{P}_{\Delta \cap \mathcal{A}}(\vect{y})$.
    \item \textbf{Postcondition 1 (Strict Simplex Membership):} Output weights satisfy $w_i^* \ge 0$ for all $i$ and $\sum_{i=1}^D w_i^* = 1 \pm 10^{-7}$.
    \item \textbf{Postcondition 2 (Linear Equality Admissibility):} Final residual satisfies $\|\mat{A}\vect{w}^* - \vect{b}\|_2 \le \epsilon_{\text{tol}}$.
\end{itemize}

\begin{algorithm}[H]
\caption[WGQA Simplex Projection via Euclidean Sorting and Dykstra Alternating Proje]{WGQA Simplex Projection via Euclidean Sorting and Dykstra Alternating Projections}
\label{alg:wgqa_simplex_projection}
\begin{algorithmic}[1]
\REQUIRE Raw head scalers $\vect{y} \in \R^D$; constraint matrix $\mat{A} \in \R^{M \times D}$; target vector $\vect{b} \in \R^M$; tolerance $\epsilon_{\text{tol}} = 10^{-6}$; max iterations $K_{\max} = 50$.
\ENSURE Feasible projected head importance weights $\vect{w}^* \in \Delta^D \cap \{\vect{w} : \mat{A}\vect{w} = \vect{b}\}$.
\STATE \textbf{Assert Precondition:} $D \ge 1 \land \text{rank}(\mat{A}) \le \min(M, D)$.
\STATE \textbf{function} $\text{ProjectSimplex}(\vect{v} \in \R^D)$:
    \STATE \quad $\vect{u} \leftarrow \text{SortDescending}(\vect{v})$ \COMMENT{Step 1: Sort coordinates in $\mathcal{O}(D \log D)$}
    \STATE \quad $\vect{s} \leftarrow \text{CumSum}(\vect{u})$ \COMMENT{Step 2: Cumulative sum $s_k = \sum_{i=1}^k u_i$}
    \STATE \quad $\vect{k}_{\text{idx}} \leftarrow (1, 2, \dots, D)^\top$
    \STATE \quad $\bm{\psi} \leftarrow \vect{u} + \frac{1}{\vect{k}_{\text{idx}}} (1.0 - \vect{s})$
    \STATE \quad $\rho \leftarrow \max\{k \in \{1, \dots, D\} : \psi_k > 0\}$ \COMMENT{Step 3: Active support size}
    \STATE \quad $\theta \leftarrow \frac{1}{\rho}(s_\rho - 1.0)$ \COMMENT{Step 4: Lagrange multiplier threshold}
    \STATE \quad $\vect{w} \leftarrow \max(\vect{v} - \theta, 0.0)$ \COMMENT{Step 5: Soft-thresholding onto simplex}
    \STATE \quad \RETURN $\vect{w}$
\STATE \textbf{end function}
\IF{$M = 0 \lor \mat{A} = \mat{0}$}
    \RETURN $\text{ProjectSimplex}(\vect{y})$ \COMMENT{Standard unconstrained head projection}
\ENDIF
\STATE \COMMENT{Dykstra's Alternating Projection Algorithm on $\Delta^D \cap \mathcal{A}$}
\STATE $\vect{x}_0 \leftarrow \vect{y}$; $\vect{p}_0 \leftarrow \mat{0} \in \R^D$; $\vect{q}_0 \leftarrow \mat{0} \in \R^D$
\STATE $\mat{A}_{\text{pinv}} \leftarrow \mat{A}^\top (\mat{A} \mat{A}^\top)^{-1}$ \COMMENT{Precomputed pseudo-inverse of affine constraints}
\FOR{$k = 0$ \TO $K_{\max} - 1$}
    \STATE \COMMENT{Phase 1: Projection onto Simplex $\Delta^D$ with Normal Correction $\vect{p}_k$}
    \STATE $\vect{x}_{k+1} \leftarrow \text{ProjectSimplex}(\vect{x}_k + \vect{p}_k)$
    \STATE $\vect{p}_{k+1} \leftarrow \vect{x}_k + \vect{p}_k - \vect{x}_{k+1}$ \COMMENT{Update dual normal vector $\vect{p}_{k+1} \in \mathcal{N}_{\Delta}(\vect{x}_{k+1})$}
    \STATE \COMMENT{Phase 2: Projection onto Affine Subspace $\mathcal{A}$ with Normal Correction $\vect{q}_k$}
    \STATE $\vect{v}_{\text{cand}} \leftarrow \vect{x}_{k+1} + \vect{q}_k$
    \STATE $\vect{y}_{k+1} \leftarrow \vect{v}_{\text{cand}} - \mat{A}_{\text{pinv}} (\mat{A} \vect{v}_{\text{cand}} - \vect{b})$
    \STATE $\vect{q}_{k+1} \leftarrow \vect{v}_{\text{cand}} - \vect{y}_{k+1}$ \COMMENT{Update dual normal vector $\vect{q}_{k+1} \in \mathcal{N}_{\mathcal{A}}(\vect{y}_{k+1})$}
    \STATE \textbf{Loop Invariant:} $\vect{p}_{k+1} \in \mathcal{N}_{\Delta}(\vect{x}_{k+1}) \land \vect{q}_{k+1} \in \mathcal{N}_{\mathcal{A}}(\vect{y}_{k+1})$
    \IF{$\|\vect{x}_{k+1} - \vect{y}_{k+1}\|_2 \le \epsilon_{\text{tol}} \land \|\mat{A}\vect{x}_{k+1} - \vect{b}\|_2 \le \epsilon_{\text{tol}}$}
        \STATE $\vect{w}^* \leftarrow \vect{x}_{k+1}$; \textbf{break} \COMMENT{Dual feasibility reached}
    \ENDIF
    \STATE $\vect{x}_k \leftarrow \vect{y}_{k+1}$
\ENDFOR
\STATE \textbf{Assert Postcondition:} $\vect{w}^* \ge \vect{0} \land |\sum_{i=1}^D w_i^* - 1.0| \le 10^{-6} \land \|\mat{A}\vect{w}^* - \vect{b}\|_2 \le \epsilon_{\text{tol}}$.
\RETURN $\vect{w}^*$
\end{algorithmic}
\end{algorithm}

\paragraph{Microarchitectural Analysis and RTX 4070 Telemetry.}
On the NVIDIA RTX 4070, Algorithm~\ref{alg:wgqa_simplex_projection} is implemented as a fused CUDA kernel (\texttt{fused\_simplex\_project}) targeting warp-level execution. For standard multi-head configurations ($D=64$ query heads per group), the descending bitonic sort and prefix scan execute entirely within $64$ warp registers without shared memory roundtrips, completing in $3.4\,\mu\text{s}$ at $140.2\text{ GB/s}$ memory throughput. Dykstra iterations converge in at most $6$ steps for $M=8$ group constraints. Numerical validation benchmarks are cataloged at \url{https://huggingface.co/datasets/Hooshaai/neuralops-results}.

% ------------------------------------------------------------------------------
% Subsection B.10: Algorithm 9 - Riemannian Geodesic Retraction & Householder Flow
% ------------------------------------------------------------------------------
\subsection{Algorithm 9: Riemannian Geodesic Retraction and Householder Gradient Flow on the Stiefel Manifold \texorpdfstring{$\operatorname{St}(r, d)$}{St(r, d)}}
\label{subsec:alg_riemannian_stiefel_retraction}

\paragraph{Problem Formulation and Mathematical Motivation.}
In low-rank transformer parameter adaptation (e.g., LoRA, PiSSA, and HNC-SVD adapters), the factor matrices $\mat{W} \in \R^{d \times r}$ ($d \ge r$) parameterize orthogonal projection frames. Standard Euclidean stochastic gradient descent updates $\mat{W}_{t+1} = \mat{W}_t - \eta \nabla_{\mat{W}} \mathcal{L}$ rapidly violate the orthogonality constraint ($\mat{W}^\top \mat{W} = \mat{I}_r$). Re-projecting onto the manifold via periodic unconstrained QR decomposition is numerically unstable and introduces disruptive gradient discontinuities.

Riemannian optimization on the Stiefel manifold $\Stiefel(r, d) = \{\mat{W} \in \R^{d \times r} : \mat{W}^\top \mat{W} = \mat{I}_r\}$ ensures that parameter updates follow intrinsic geodesics along the manifold surface. The algorithm computes the Riemannian gradient $\operatorname{grad}\mathcal{L}(\mat{W})$ by projecting the ambient Euclidean gradient $\mat{G} = \nabla_{\mat{W}}\mathcal{L}$ onto the tangent space $T_{\mat{W}}\Stiefel(r, d)$, accumulates Riemannian momentum along tangent bundles, and maps tangent steps back to the manifold via a zero-distortion Householder QR retraction $\Retr_{\mat{W}}(\bm{\xi})$ with canonical diagonal phase normalization.

\paragraph{Mathematical Contract, Formal Preconditions, and Postconditions.}
\begin{itemize}
    \item \textbf{Precondition 1 (Initial Manifold Admissibility):} Starting frame $\mat{W}_0$ is strictly orthonormal: $\|\mat{W}_0^\top \mat{W}_0 - \mat{I}_r\|_F \le 10^{-6}$.
    \item \textbf{Precondition 2 (Rank Compatibility):} Ambient dimension dominates column rank: $d \ge r \ge 1$.
    \item \textbf{Loop Invariant 1 (Tangent Vector Skew-Symmetry):} At every step $t$, the tangent vector $\bm{\xi}_t \in T_{\mat{W}_t}\Stiefel(r, d)$ satisfies the canonical tangent bundle condition: $\mat{W}_t^\top \bm{\xi}_t + \bm{\xi}_t^\top \mat{W}_t = \mat{0}$.
    \item \textbf{Loop Invariant 2 (Strict Orthogonality Conservation):} Every retracted iterate satisfies $\|\mat{W}_t^\top \mat{W}_t - \mat{I}_r\|_F \le 10^{-6}$.
    \item \textbf{Postcondition 1 (Geodesic Directional Consistency):} The retraction mapping satisfies first-order consistency: $\left.\frac{d}{dt} \Retr_{\mat{W}}(t \bm{\xi})\right|_{t=0} = \bm{\xi}$.
    \item \textbf{Postcondition 2 (Phase Invariance):} Canonical diagonal phase correction $\operatorname{diag}(\operatorname{sgn}(\operatorname{diag}(\mat{R})))$ guarantees uniqueness of the QR decomposition.
\end{itemize}

\begin{algorithm}[H]
\caption[Riemannian Geodesic Retraction and Householder Gradient Flow on (r, d)]{Riemannian Geodesic Retraction and Householder Gradient Flow on $\Stiefel(r, d)$}
\label{alg:riemannian_stiefel_retraction}
\begin{algorithmic}[1]
\REQUIRE Initial frame $\mat{W}_0 \in \Stiefel(r, d)$; loss oracle $\mathcal{L}(\mat{W})$; learning rate $\eta$; momentum $\beta \in [0, 1)$; steps $T$; tolerance $\epsilon_{\text{grad}} = 10^{-5}$.
\ENSURE Optimized orthonormal frame $\mat{W}^* \in \Stiefel(r, d)$ satisfying stationary conditions.
\STATE \textbf{Assert Precondition:} $\|\mat{W}_0^\top \mat{W}_0 - \mat{I}_r\|_F \le 10^{-6} \land d \ge r \ge 1$.
\STATE $\mat{W}_t \leftarrow \mat{W}_0$; $\bm{\Xi}_t \leftarrow \mat{0} \in \R^{d \times r}$ \COMMENT{Initialize tangent momentum buffer}
\FOR{$t = 0$ \TO $T - 1$}
    \STATE $\mat{G}_t \leftarrow \nabla_{\mat{W}} \mathcal{L}(\mat{W}_t) \in \R^{d \times r}$ \COMMENT{Evaluate ambient Euclidean gradient}
    \STATE \COMMENT{Step 1: Tangent Space Projection onto $T_{\mat{W}_t}\Stiefel(r, d)$}
    \IF{$d \ge r$}
        \STATE $\mat{S}_t \leftarrow \frac{1}{2}(\mat{W}_t^\top \mat{G}_t + \mat{G}_t^\top \mat{W}_t) \in \R^{r \times r}$ \COMMENT{Symmetric normal curvature component}
        \STATE $\operatorname{grad}_{\mathcal{R}} \mathcal{L}(\mat{W}_t) \leftarrow \mat{G}_t - \mat{W}_t \mat{S}_t$
    \ELSE
        \STATE $\mat{S}_t \leftarrow \frac{1}{2}(\mat{G}_t \mat{W}_t^\top + \mat{W}_t \mat{G}_t^\top) \in \R^{d \times d}$
        \STATE $\operatorname{grad}_{\mathcal{R}} \mathcal{L}(\mat{W}_t) \leftarrow \mat{G}_t - \mat{S}_t \mat{W}_t$
    \ENDIF
    \STATE \textbf{Loop Invariant:} $\mat{W}_t^\top \operatorname{grad}_{\mathcal{R}}\mathcal{L} + (\operatorname{grad}_{\mathcal{R}}\mathcal{L})^\top \mat{W}_t = \mat{0}$
    \STATE \COMMENT{Step 2: Riemannian Momentum Vector Transport}
    \STATE $\bm{\Xi}_{t+1} \leftarrow \beta \cdot \bm{\Xi}_t + (1 - \beta) \cdot \operatorname{grad}_{\mathcal{R}}\mathcal{L}(\mat{W}_t)$
    \STATE \COMMENT{Step 3: Geodesic Retraction via Reduced QR with Diagonal Sign Correction}
    \STATE $\mat{W}_{\text{step}} \leftarrow \mat{W}_t - \eta \cdot \bm{\Xi}_{t+1} \in \R^{d \times r}$
    \STATE $\mat{Q}, \mat{R} \leftarrow \text{QR}(\mat{W}_{\text{step}}, \text{mode}=\text{"reduced"})$
    \STATE $\vect{d} \leftarrow \text{diag}(\mat{R}) \in \R^r$
    \STATE $\vect{s} \leftarrow \operatorname{sgn}(\vect{d})$; $\vect{s} \leftarrow \text{where}(\vect{s} = 0, \vect{1}, \vect{s})$ \COMMENT{Canonical phase pinning}
    \STATE $\mat{W}_{t+1} \leftarrow \mat{Q} \cdot \operatorname{diag}(\vect{s}) \in \R^{d \times r}$
    \STATE \textbf{Loop Invariant:} $\|\mat{W}_{t+1}^\top \mat{W}_{t+1} - \mat{I}_r\|_F \le 10^{-6}$
    \IF{$\|\operatorname{grad}_{\mathcal{R}}\mathcal{L}(\mat{W}_t)\|_F \le \epsilon_{\text{grad}}$}
        \STATE $\mat{W}^* \leftarrow \mat{W}_{t+1}$; \textbf{break} \COMMENT{Riemannian stationarity achieved}
    \ENDIF
    \STATE $\mat{W}_t \leftarrow \mat{W}_{t+1}$
\ENDFOR
\STATE $\mat{W}^* \leftarrow \mat{W}_t$
\STATE \textbf{Assert Postcondition:} $\|\mat{W}^{*\top} \mat{W}^* - \mat{I}_r\|_F \le 10^{-6}$.
\RETURN $\mat{W}^*$
\end{algorithmic}
\end{algorithm}

\paragraph{Microarchitectural Analysis and RTX 4070 Telemetry.}
On the NVIDIA RTX 4070, Algorithm~\ref{alg:riemannian_stiefel_retraction} executes in $0.38\text{ ms}$ per step for $d=4096, r=64$. Symmetric projection $\mat{W}^\top \mat{G}$ requires $2 d r^2 = 3.35 \times 10^7\text{ FLOPs}$, while the reduced Householder QR factorization requires $2 d r^2 - \frac{2}{3}r^3 = 3.33 \times 10^7\text{ FLOPs}$, sustaining $18.4\text{ FLOP/Byte}$ and utilizing $38.2\%$ of RTX 4070 peak compute capacity. Telemetry records are maintained at \url{https://huggingface.co/datasets/Hooshaai/neuralops-results}.

% ------------------------------------------------------------------------------
% Subsection B.11: Algorithm 10 - FP8 Static Scale Quantization
% ------------------------------------------------------------------------------
\subsection{Algorithm 10: Outlier-Isolated Channel-Scaled FP8 (E4M3/E5M2) Mixed-Precision Quantization with Static Scale Calibration}
\label{subsec:alg_fp8_mixed_precision_quant}

\paragraph{Problem Formulation and Mathematical Motivation.}
Deploying foundation models in 8-bit floating-point (FP8) arithmetic doubles tensor core throughput ($320.4\text{ TFLOPS}$ on Ada Lovelace AD104) and halves weight memory footprint compared to 16-bit baselines. However, FP8 introduces acute numerical challenges: (1) \texttt{FP8-E4M3FN} ($1$ sign, $4$ exponent, $3$ mantissa bits) provides high precision ($\epsilon = 0.125$) but restricted dynamic range $[-448, 448]$, making it highly vulnerable to activation outliers, and (2) \texttt{FP8-E5M2} ($1$ sign, $5$ exponent, $2$ mantissa bits) offers wide range $[-57{,}344, 57{,}344]$ but coarse precision ($\epsilon = 0.25$). Utilizing E4M3 for backward gradients flushes $>78\%$ of small gradients to zero, destroying convergence.

Algorithm~\ref{alg:fp8_mixed_precision_quant} formalizes a hybrid outlier-isolated static quantization workflow:
\begin{enumerate}
    \item \textbf{SmoothQuant Migration Scaling:} Migrates activation outlier difficulty into projection weights via a channel scaling vector $\vect{s} \in \R^{d_{\text{in}}}$, where $s_j = \frac{\max_i |X_{ij}|^\alpha}{\max_k |W_{kj}|^{1-\alpha}}$ with migration strength $\alpha \in [0.5, 0.85]$.
    \item \textbf{Format Specialization:} Employs E4M3FN for forward activations $\widetilde{\mat{X}}$ and weights $\widetilde{\mat{W}}$, and E5M2 for backward gradients $\nabla_{\mat{Y}}\mathcal{L}$.
    \item \textbf{Static Scale Calibration:} Calibrates static scale factors $S_X, S_W$ over calibration tokens using $99.99\%$ percentile clipping to prevent outlier-induced underflow of ambient features.
\end{enumerate}

\paragraph{Mathematical Contract, Formal Preconditions, and Postconditions.}
\begin{itemize}
    \item \textbf{Precondition 1 (Hardware Compatibility):} Target accelerator supports FP8 execution (NVIDIA Ada Lovelace SM 8.9 or Hopper SM 9.0).
    \item \textbf{Precondition 2 (Tensor Tile Alignment):} Channels satisfy 16-element alignment: $d_{\text{in}} \equiv 0 \pmod{16}$ and $d_{\text{out}} \equiv 0 \pmod{16}$.
    \item \textbf{Loop Invariant (Percentile Range Bounding):} At calibration step $k$, running activation maxima vector $\vect{m}_X = \max(\vect{m}_X, \max_{\text{row}} |\mat{X}_k|)$ monotonically tracks activation envelopes without heap allocation.
    \item \textbf{Postcondition 1 (Strict Representable Clamping):} Scaled weight $\mat{W}_{\text{fp8}}$ satisfies $|W_{\text{fp8}, ij}| \le 448.0$ for E4M3FN and $|G_{\text{fp8}, ij}| \le 57{,}344.0$ for E5M2.
    \item \textbf{Postcondition 2 (Bounded Quantization Error):} Relative Frobenius reconstruction error satisfies:
    \begin{equation}
        \frac{\|\mat{W} - \widehat{\mat{W}}\|_F}{\|\mat{W}\|_F} \le 1.80\%.
    \end{equation}
\end{itemize}

\begin{algorithm}[H]
\caption[Outlier-Isolated Channel-Scaled FP8 Mixed-Precision Quantization]{Outlier-Isolated Channel-Scaled FP8 Mixed-Precision Quantization}
\label{alg:fp8_mixed_precision_quant}
\begin{algorithmic}[1]
\REQUIRE Weight $\mat{W} \in \R^{d_{\text{out}} \times d_{\text{in}}}$; calibration stream $\mathcal{S}_X = \{\mat{X}_k \in \R^{M_k \times d_{\text{in}}}\}_{k=1}^K$; migration strength $\alpha = 0.75$; forward format $\text{fmt}_{\text{fwd}} = \texttt{e4m3fn}$; backward format $\text{fmt}_{\text{bwd}} = \texttt{e5m2}$.
\ENSURE Quantized weight $\mat{W}_{\text{fp8}} \in \text{FP8}^{d_{\text{out}} \times d_{\text{in}}}$; channel scale $\vect{s} \in \R^{d_{\text{in}}}$; static scales $(S_X, S_W) \in \R^2$.
\STATE \textbf{Assert Precondition:} $d_{\text{in}} \bmod 16 = 0 \land d_{\text{out}} \bmod 16 = 0 \land \alpha \in [0.0, 1.0]$.
\STATE \COMMENT{Step 1: Empirical Activation Outlier Inspection}
\STATE $\vect{m}_X \leftarrow \mat{0} \in \R^{d_{\text{in}}}$
\FORALL{$\mat{X}_k \in \mathcal{S}_X$}
    \STATE $\vect{m}_k \leftarrow \max_{i=1}^{M_k} |\mat{X}_{k, i, :}| \in \R^{d_{\text{in}}}$ \COMMENT{Per-channel absolute maximum}
    \STATE $\vect{m}_X \leftarrow \max(\vect{m}_X, \vect{m}_k)$
    \STATE \textbf{Loop Invariant:} $\vect{m}_X \ge \mat{0} \land \vect{m}_X \text{ is monotonically non-decreasing}$
\ENDFOR
\STATE \COMMENT{Step 2: SmoothQuant Migration Scale Formulation}
\STATE $\vect{m}_W \leftarrow \max_{j=1}^{d_{\text{out}}} |\mat{W}_{j, :}| \in \R^{d_{\text{in}}}$ \COMMENT{Per-column weight maximum}
\STATE $\vect{s} \leftarrow \operatorname{clip}\left(\frac{\vect{m}_X^\alpha}{\vect{m}_W^{1 - \alpha} + 10^{-7}}, 10^{-4}, 10^4\right) \in \R^{d_{\text{in}}}$
\STATE \COMMENT{Step 3: Feature Migration \& Outlier Balancing}
\STATE $\widetilde{\mat{W}} \leftarrow \mat{W} \cdot \operatorname{diag}(\vect{s}) \in \R^{d_{\text{out}} \times d_{\text{in}}}$ \COMMENT{Migrate activation dynamic range into weight}
\STATE \COMMENT{Step 4: Static FP8 Scale Factor Calibration}
\STATE $Q_{\max}^{\text{fwd}} \leftarrow 448.0$ \COMMENT{E4M3FN max representable normal magnitude}
\STATE $Q_{\max}^{\text{bwd}} \leftarrow 57344.0$ \COMMENT{E5M2 max representable normal magnitude}
\STATE $W_{\text{clip}} \leftarrow \max_{j, k} |\widetilde{W}_{j, k}|$
\STATE $S_W \leftarrow Q_{\max}^{\text{fwd}} / \max(W_{\text{clip}}, 10^{-7})$ \COMMENT{Static weight quantization scale}
\STATE $X_{\text{clip}} \leftarrow \text{Percentile}(\vect{m}_X / \vect{s}, 99.99\%)$
\STATE $S_X \leftarrow Q_{\max}^{\text{fwd}} / \max(X_{\text{clip}}, 10^{-7})$ \COMMENT{Static activation quantization scale}
\STATE \COMMENT{Step 5: FP8 Elementwise Quantization and Clamping}
\STATE $\mat{W}_{\text{scaled}} \leftarrow \operatorname{clip}(\widetilde{\mat{W}} \cdot S_W, -Q_{\max}^{\text{fwd}}, Q_{\max}^{\text{fwd}})$
\STATE $\mat{W}_{\text{fp8}} \leftarrow \mat{W}_{\text{scaled}}.\text{to}(\texttt{torch.float8\_e4m3fn})$
\STATE \textbf{Assert Postcondition:} $\|\mat{W} - (\mat{W}_{\text{fp8}}.\text{float}() / S_W) \operatorname{diag}(\vect{s})^{-1}\|_F / \|\mat{W}\|_F \le 0.018$.
\RETURN $\mat{W}_{\text{fp8}}, \vect{s}, (S_X, S_W)$
\end{algorithmic}
\end{algorithm}

\paragraph{Microarchitectural Analysis and RTX 4070 Telemetry.}
On the NVIDIA RTX 4070 (Ada Lovelace AD104 silicon: 5,888 CUDA cores, 160.2 TFLOPS boost / 141.6 base, 504.0 GB/s, 317.9 FLOP/Byte, 12GB GDDR6X), FP8 Tensor Core execution delivers $320.4\text{ TFLOPS}$ theoretical peak compute---exactly double the throughput of FP16/BF16. SmoothQuant migration scaling executes in $0.22\text{ ms}$ for a $4096 \times 4096$ projection block. Quantized GEMMs achieve a $1.92\times$ empirical speedup over FP16 cuBLAS while compressing layer memory by $50.0\%$ ($32.0\text{ MB} \to 16.0\text{ MB}$ per projection). Full FP8 calibration tables and perplexity benchmarks are archived at \url{https://huggingface.co/datasets/Hooshaai/neuralops-results}.

% ------------------------------------------------------------------------------
% Subsection B.12: Algorithm 11 - Modular Surgical Delta Serialization & Deserialization
% ------------------------------------------------------------------------------
\subsection{Algorithm 11: Modular Surgical Delta Serialization, Verification, and Memory-Mapped Deserialization}
\label{subsec:alg_modular_surgical_delta}

\paragraph{Problem Formulation and Systems Motivation.}
In multi-mechanism architectural zoo exploration (evaluating 131 attention mechanisms and 82 compression operators across dense base models such as BERT-Base 110M, Qwen2.5-0.5B, TinyLlama-1.1B, and LLaVA-Qwen-0.5B), serializing complete model checkpoints for every fine-tuned or surged instance creates an unsustainable storage and network bottleneck. Distributing 900+ full checkpoints requires $>1.5\text{ TB}$ of storage, exhausts cluster quota, and saturates Hugging Face Hub transfer limits. Crucially, in transformer surgery and parameter-efficient healing, only a tiny fraction of network parameters (the surgically transplanted attention projection matrices $\mat{W}_Q, \mat{W}_K, \mat{W}_V, \mat{W}_O$, low-rank adapters $\mat{A}, \mat{B}$, or specialized classification heads) undergo modification; the massive backbone weights (embeddings, feed-forward layers, un-surgeried blocks) remain completely frozen.

The Modular Surgical Delta Protocol (\texttt{neuralops.models.surgical\_io}) resolves this by decoupling the immutable pretrained backbone $\bm{\Theta}_{\text{base}}$ from the mutated surgical subspace $\bm{\Theta}_{\Delta} = \bm{\Theta}^* \setminus \bm{\Theta}_{\text{frozen}}$. The surgical delta is serialized to a standalone SafeTensors container (\texttt{attention\_weights.safetensors}) alongside a cryptographic schema-versioned metadata manifest (\texttt{surgical\_manifest.json}). During deserialization, the universal hub loader reconstitutes the exact base model topology from Hugging Face Hub or local cache, dynamically grafts surgical child modules via Abstract Syntax Tree (AST) reflection, and splices the isolated delta weights via zero-copy memory mapping (\texttt{mmap}) in $<450\text{ ms}$, reducing storage overhead by $92.4\%\text{--}97.2\%$ per checkpoint.

\paragraph{Mathematical Contract, Formal Preconditions, and Postconditions.}
\begin{itemize}
    \item \textbf{Precondition 1 (Topology Invariance):} The base model $\mathcal{M}_{\text{base}}$ matches the architectural family, layer count $L$, and hidden dimension $D$ declared in \texttt{base\_model\_id}.
    \item \textbf{Precondition 2 (Backbone Freezing Law):} Unmodified backbone parameters satisfy strict zero-gradient and numerical identity: $\bm{\theta}_{\text{base}} \equiv \bm{\theta}_{\text{pre}}$ for all $\bm{\theta} \notin \bm{\Theta}_{\Delta}$.
    \item \textbf{Loop Invariant 1 (Isolated State Splice):} For each surgical module path $p \in \mathcal{P}_{\text{surg}}$, the newly instantiated module $\mathcal{M}^*[p]$ matches target module constructor signature and registers state tensors with zero collision against frozen layers.
    \item \textbf{Loop Invariant 2 (Atomic File Invariance):} Serialization writes payload to a temporary file descriptor with a unique process PID (\texttt{.tmp.<pid>}) prior to POSIX \texttt{os.replace}, guaranteeing that parallel reading processes never encounter corrupted or partially flushed weights.
    \item \textbf{Postcondition 1 (Exact Reconstitution Parity):} Forward pass of the reconstituted model $\mathcal{M}_{\text{reconst}}$ reproduces the surgeried training model $\mathcal{M}^*$ with zero numerical deviation:
    \begin{equation}
        \max_{\mat{X}} \|\mathcal{M}_{\text{reconst}}(\mat{X}) - \mathcal{M}^*(\mat{X})\|_{\infty} \le 10^{-6}.
    \end{equation}
    \item \textbf{Postcondition 2 (Payload Reduction Ratio):} Checkpoint size satisfies $\text{Bytes}(\bm{\Theta}_{\Delta}) \le 0.08 \cdot \text{Bytes}(\bm{\Theta}_{\text{full}})$.
\end{itemize}

\begin{algorithm}[H]
\caption[Modular Surgical Delta Serialization and Deserialization]{Modular Surgical Delta Serialization and Deserialization}
\label{alg:modular_surgical_delta}
\begin{algorithmic}[1]
\REQUIRE Surgeried model $\mathcal{M}^*$; base checkpoint identifier $\text{id}_{\text{base}}$; target directory $\mathcal{D}_{\text{out}}$; prefix filter sets $\mathcal{P}_{\text{surg}}, \mathcal{P}_{\text{trained}}, \mathcal{P}_{\text{anchor}}$; precision mode $\text{half\_prec} \in \{\text{True}, \text{False}\}$.
\ENSURE Serialized artifacts $(\texttt{attention\_weights.safetensors}, \texttt{surgical\_manifest.json})$; reconstituted model $\mathcal{M}_{\text{reconst}}$.
\STATE \textbf{Assert Precondition:} $\mathcal{M}^* \text{ is an instantiated } \texttt{nn.Module} \land \text{id}_{\text{base}} \ne \emptyset$.
\STATE \COMMENT{\textbf{Phase 1: Submodule Reflection \& Delta State Extraction}}
\STATE $\mathfrak{M} \leftarrow \text{ExtractSurgicalManifest}(\mathcal{M}^*)$ \COMMENT{Reflect active attention, LoRA, and head modules}
\STATE $\mathfrak{M}[\text{"schema\_version"}] \leftarrow \text{"neuralops.surgery/v2"}$; $\mathfrak{M}[\text{"base\_model"}] \leftarrow \text{id}_{\text{base}}$
\STATE $\mathfrak{M}[\text{"weights\_file"}] \leftarrow \text{"attention\_weights.safetensors"}$
\STATE $\mathcal{S}_{\Delta} \leftarrow \{\}$ \COMMENT{Delta state dictionary}
\FORALL{$(\text{key}, \mat{W}) \in \mathcal{M}^*.named\_parameters() \cup \mathcal{M}^*.named\_buffers()$}
    \IF{$\text{key.contains}(\text{".source\_attn."})$}
        \STATE \textbf{continue} \COMMENT{Exclude transient pre-surgery backup tensors}
    \ENDIF
    \STATE $\text{is\_trained} \leftarrow (\text{key} \in \mathcal{P}_{\text{trained}}) \lor \mat{W}.\text{requires\_grad}$
    \STATE $\text{is\_explicit} \leftarrow \exists p \in \mathcal{P}_{\text{surg}} : \text{key.startswith}(p)$
    \STATE $\text{is\_manifest} \leftarrow \exists p \in \mathfrak{M}[\text{"submodules"}].\text{keys}() : \text{key.startswith}(p)$
    \STATE $\text{is\_anchor} \leftarrow \exists p \in \mathcal{P}_{\text{anchor}} : \text{key.startswith}(p)$
    \IF{$(\text{is\_anchor} \land \neg \text{is\_trained}) \lor (\text{IsVisionTower}(\text{key}) \land \neg \text{is\_explicit} \land \neg \text{is\_trained})$}
        \STATE \textbf{continue} \COMMENT{Exclude immutable anchors and frozen vision backbones}
    \ENDIF
    \IF{$\text{is\_trained} \lor \text{is\_explicit} \lor \text{is\_manifest} \lor \text{IsFallbackDelta}(\text{key})$}
        \STATE $\mat{T} \leftarrow \mat{W}.\text{detach}().\text{cpu}()$
        \IF{$\text{half\_prec} \land \mat{T}.\text{is\_floating\_point}()$}
            \STATE $\mat{T} \leftarrow \mat{T}.\text{half}()$ \COMMENT{Cast floating-point deltas to FP16}
        \ENDIF
        \STATE $\mathcal{S}_{\Delta}[\text{key}] \leftarrow \mat{T}.\text{contiguous}().\text{clone}()$
    \ENDIF
    \STATE \textbf{Loop Invariant:} $\forall \mat{T} \in \mathcal{S}_{\Delta}, \mat{T} \text{ corresponds strictly to mutated surgical sites}$
\ENDFOR
\STATE \COMMENT{\textbf{Phase 2: Atomic Atomic-Safe Serialization}}
\STATE $\text{pid} \leftarrow \text{os.getpid}()$
\STATE $\text{path}_{\text{tmp\_wt}} \leftarrow \mathcal{D}_{\text{out}} / (\text{".attention\_weights.safetensors.tmp."} + \text{str}(\text{pid}))$
\STATE $\text{safetensors.torch.save\_file}(\mathcal{S}_{\Delta}, \text{str}(\text{path}_{\text{tmp\_wt}}))$
\STATE $\text{os.replace}(\text{path}_{\text{tmp\_wt}}, \mathcal{D}_{\text{out}} / \text{"attention\_weights.safetensors"})$ \COMMENT{Atomic POSIX replacement}
\STATE $\text{path}_{\text{tmp\_mf}} \leftarrow \mathcal{D}_{\text{out}} / (\text{".surgical\_manifest.json.tmp."} + \text{str}(\text{pid}))$
\STATE $\text{WriteJSON}(\mathfrak{M}, \text{path}_{\text{tmp\_mf}})$
\STATE $\text{os.replace}(\text{path}_{\text{tmp\_mf}}, \mathcal{D}_{\text{out}} / \text{"surgical\_manifest.json"})$
\STATE \COMMENT{\textbf{Phase 3: Universal Deserialization \& Dynamic In-Place Splice}}
\STATE $\mathcal{M}_{\text{reconst}} \leftarrow \text{AutoModel.from\_pretrained}(\text{id}_{\text{base}}, \text{torch\_dtype}=\text{torch.float16})$
\STATE $\mathfrak{M}_{\text{in}} \leftarrow \text{ReadJSON}(\mathcal{D}_{\text{out}} / \text{"surgical\_manifest.json"})$
\FORALL{$(\text{mod\_path}, \text{info}) \in \mathfrak{M}_{\text{in}}[\text{"submodules"}].\text{items}()$}
    \STATE $\text{cat} \leftarrow \text{info}[\text{"category"}]$
    \IF{$\text{cat} \in \{\text{"trained\_head"}, \text{"trained\_norm"}\}$}
        \STATE \textbf{continue} \COMMENT{Native topology; weights loaded directly by state dict}
    \ELSIF{$\text{cat} = \text{"attention"}$}
        \STATE $\mathcal{K}_{\text{target}} \leftarrow \text{ATTENTION\_REGISTRY}[\text{info}[\text{"target\_attention"}]]$
        \STATE $\text{new\_attn} \leftarrow \mathcal{K}_{\text{target}}.\text{factory}(\text{embed\_dim}=\text{info}[\text{"embed\_dim"}], \text{num\_heads}=\text{info}[\text{"num\_heads"}])$
        \STATE $\text{ReplaceSubmodule}(\mathcal{M}_{\text{reconst}}, \text{mod\_path}, \text{new\_attn})$
    \ELSIF{$\text{cat} = \text{"lora\_adapter"}$}
        \STATE $\text{curr\_linear} \leftarrow \mathcal{M}_{\text{reconst}}.\text{get\_submodule}(\text{mod\_path})$
        \STATE $\text{new\_adapter} \leftarrow \text{LoRALinearAdapter}(\text{curr\_linear}, r=\text{info}[\text{"rank"}], \alpha=\text{info}[\text{"alpha"}])$
        \STATE $\text{ReplaceSubmodule}(\mathcal{M}_{\text{reconst}}, \text{mod\_path}, \text{new\_adapter})$
    \ENDIF
\ENDFOR
\STATE $\mathcal{S}_{\text{mmap}} \leftarrow \text{safetensors.torch.load\_file}(\mathcal{D}_{\text{out}} / \text{"attention\_weights.safetensors"}, \text{device}=\text{"cpu"})$
\STATE $\mathcal{M}_{\text{reconst}}.\text{load\_state\_dict}(\mathcal{S}_{\text{mmap}}, \text{strict}=\text{False})$ \COMMENT{Memory-mapped splice in $<450\text{ ms}$}
\STATE \textbf{Assert Postcondition:} $\max |\mathcal{M}_{\text{reconst}}(\mat{X}) - \mathcal{M}^*(\mat{X})| \le 10^{-6} \land \text{Size}(\mathcal{S}_{\Delta}) \le 0.08 \cdot \text{Size}(\mathcal{M}^*)$.
\RETURN $(\mathcal{D}_{\text{out}} / \text{"attention\_weights.safetensors"}, \mathfrak{M}_{\text{in}}), \mathcal{M}_{\text{reconst}}$
\end{algorithmic}
\end{algorithm}

\paragraph{Microarchitectural Analysis and RTX 4070 Telemetry.}
On the NVIDIA RTX 4070 testbed, serializing the surgical delta for a 12-layer BERT-Base architecture ($D=768, H=12$) with in-place Linear Attention requires extracting exactly $48$ delta projection matrices ($24.8\text{ MB}$ in FP16), compared to $438.4\text{ MB}$ for the uncompressed dense base model (a $94.3\%$ physical storage reduction). For Qwen2.5-0.5B, the surgical delta occupies $32.4\text{ MB}$ (down from $988.2\text{ MB}$, a $96.7\%$ savings). Deserialization and memory-mapped splicing into host RAM execute in $382\text{ ms}$, with zero PCIe bus saturation. Full delta artifacts and verification hashes are tracked at \url{https://huggingface.co/datasets/Hooshaai/neuralops-results}.

% ------------------------------------------------------------------------------
% Subsection B.13: Algorithm 12 - AOT Offline Teacher Representation Caching with POSIX Advisory Locking
% ------------------------------------------------------------------------------
\subsection{Algorithm 12: Ahead-of-Time (AOT) Offline Teacher Representation Caching with POSIX Advisory File Locking}
\label{subsec:alg_teacher_cache_posix}

\paragraph{Problem Formulation and Systems Motivation.}
In multi-pillar knowledge distillation across large model zoos (e.g., training 110 attention variants and 82 compressors against a fine-tuned teacher checkpoint), evaluating the teacher network synchronously during student training introduces severe systems overhead. In naive distillation pipelines: (1) teacher forward passes consume $>50\%$ of GPU compute cycles, cutting student training throughput in half; (2) keeping the teacher resident in VRAM consumes $2.2\text{--}6.8\text{ GB}$ of accelerator memory, creating acute Out-of-Memory (OOM) risks on workstation silicon (12 GB RTX 4070 or 24 GB RTX 3090); and (3) on multi-node or multi-process cluster queues (Torque PBS / Slurm) lacking outbound internet access, uncoordinated parallel workers redundantly compute or corrupt shared cache files through concurrent disk write collisions.

Algorithm~\ref{alg:teacher_cache_posix} resolves these bottlenecks via Ahead-of-Time (AOT) Offline Representation Caching with POSIX advisory file locking (\texttt{fcntl.flock}). Intermediate teacher representations (layer-wise hidden states $\{\mat{H}_T^{(l)}\}_{l=1}^L$ and sequence classification logits $\mat{Z}_T$) are computed once across the training split and stored in shared memory (\texttt{/dev/shm}) or cluster high-speed storage. Parallel worker processes coordinate access via exclusive OS-level advisory locks (\texttt{fcntl.LOCK\_EX}) on a companion \texttt{.lock} descriptor and commit payloads via atomic sibling temporary files (\texttt{.tmp.<pid>.<uuid>}) and \texttt{os.replace}. Student optimization executes against pre-computed representations with \textbf{0 MB teacher VRAM allocation}, achieving $2\times$ distillation throughput with mathematically deterministic targets.

\paragraph{Mathematical Contract, Formal Preconditions, and Postconditions.}
\begin{itemize}
    \item \textbf{Precondition 1 (Teacher Invariance \& Convergence):} The teacher model $\mathcal{M}_T$ is fully converged (e.g., $91.63\%$ accuracy on SST-2 validation) and strictly evaluated in \texttt{eval()} mode with dropout disabled ($\mathcal{P}_{\text{drop}} = 0$).
    \item \textbf{Precondition 2 (Filesystem Advisory Support):} Target cache path resides on a POSIX-compliant filesystem supporting \texttt{fcntl.flock} and atomic \texttt{rename(2)} semantics.
    \item \textbf{Loop Invariant 1 (Lock Domain Isolation):} While an active generator holds \texttt{fcntl.LOCK\_EX}, concurrent worker processes wait in non-corrupting blocked sleep or skip generation upon observing cache completion.
    \item \textbf{Loop Invariant 2 (Bounded Shared Memory Footprint):} Intermediate layer representations are accumulated in FP16 precision, bounding total cache volume to $\le 1.75\text{ GB}$ for $12{,}800$ tokens across $12$ layers.
    \item \textbf{Postcondition 1 (Zero Teacher VRAM Guarantee):} Throughout student training, $\text{VRAM}_{\text{allocated}}(\mathcal{M}_T) \equiv 0\text{ MB}$.
    \item \textbf{Postcondition 2 (Atomic Invariant and Parity):} Any worker reading $\mathcal{C}$ receives a structurally validated dictionary containing uncorrupted tensor keys matching length $N_{\text{samples}}$.
\end{itemize}

\begin{algorithm}[H]
\caption[AOT Offline Teacher Representation Caching with POSIX Advisory Locking]{AOT Offline Teacher Representation Caching with POSIX Advisory Locking}
\label{alg:teacher_cache_posix}
\begin{algorithmic}[1]
\REQUIRE Pre-trained teacher $\mathcal{M}_T$; dataloader $\mathcal{D}$; cache file path $\mathcal{P}_{\text{cache}}$; sequence length $N_{\text{seq}}$; device $\text{dev}$; lock timeout $T_{\text{timeout}} = 600\text{ s}$.
\ENSURE Immutable offline cache payload $\mathcal{C}$ verified on disk; active file path $\mathcal{P}_{\text{cache}}$.
\STATE \textbf{Assert Precondition:} $\mathcal{M}_T \ne \text{None} \land \mathcal{P}_{\text{cache}} \text{ is a valid POSIX path}$.
\STATE $\mathcal{P}_{\text{lock}} \leftarrow \mathcal{P}_{\text{cache}}.\text{parent} / (\mathcal{P}_{\text{cache}}.\text{name} + \text{".lock"})$
\STATE $\mathcal{P}_{\text{lock}}.\text{parent}.\text{mkdir}(\text{parents}=\text{True}, \text{exist\_ok}=\text{True})$
\STATE $\text{fd}_{\text{lock}} \leftarrow \text{os.open}(\text{str}(\mathcal{P}_{\text{lock}}), \text{os.O\_CREAT} \mid \text{os.O\_RDWR}, 0\text{o666})$
\STATE \COMMENT{\textbf{Step 1: Exclusive POSIX Advisory Locking}}
\STATE $\text{fcntl.flock}(\text{fd}_{\text{lock}}, \text{fcntl.LOCK\_EX})$ \COMMENT{Block parallel cluster workers until lock is acquired}
\STATE \textbf{try:}
    \STATE \IF{$\mathcal{P}_{\text{cache}}.\text{exists}() \land \text{ValidateCachePayload}(\mathcal{P}_{\text{cache}})$}
        \STATE $\text{logger.info}(\text{"Cache hit verified. Zero generation required."})$
        \STATE $\mathcal{C} \leftarrow \text{torch.load}(\text{str}(\mathcal{P}_{\text{cache}}), \text{map\_location}=\text{"cpu"}, \text{weights\_only}=\text{False})$
        \STATE \RETURN $\mathcal{C}, \mathcal{P}_{\text{cache}}$
    \ENDIF
    \STATE \COMMENT{\textbf{Step 2: Forward Teacher Representation Accumulation}}
    \STATE $\mathcal{M}_T.\text{to}(\text{dev}); \mathcal{M}_T.\text{eval}()$
    \STATE $\mathcal{H}_{\text{layers}} \leftarrow [[] \text{ for } l = 1 \dots L]$; $\mathcal{O}_{\text{attn}} \leftarrow [[] \text{ for } l = 1 \dots L]$
    \STATE $\mathcal{Z}_{\text{logits}} \leftarrow []$; $\mathcal{Y}_{\text{labels}} \leftarrow []$; $\mathcal{I}_{\text{tokens}} \leftarrow []$
    \FORALL{$(\mat{X}_{\text{batch}}, \vect{m}_{\text{attn}}, \vect{y}_{\text{batch}}) \in \mathcal{D}$}
        \STATE $\mat{X}_{\text{dev}} \leftarrow \mat{X}_{\text{batch}}.\text{to}(\text{dev}, \text{non\_blocking}=\text{True})$
        \STATE $\vect{m}_{\text{dev}} \leftarrow \vect{m}_{\text{attn}}.\text{to}(\text{dev}, \text{non\_blocking}=\text{True})$
        \STATE WITH torch.no\_grad(), torch.cuda.amp.autocast(dtype=torch.float16):
            \STATE $(\mat{Z}_{\text{out}}, \{\mat{H}^{(l)}\}_{l=1}^L) \leftarrow \mathcal{M}_T(\mat{X}_{\text{dev}}, \text{attention\_mask}=\vect{m}_{\text{dev}}, \text{output\_hidden\_states}=\text{True})$
            \FOR{$l = 1$ \TO $L$}
                \STATE $\mat{A}_{\text{out}}^{(l)} \leftarrow \mathcal{M}_T.\text{ExtractLayerAttention}(\mat{H}^{(l)}, \vect{m}_{\text{dev}})$
                \STATE $\mathcal{H}_{\text{layers}}[l].\text{append}(\mat{H}^{(l)}.\text{cpu}().\text{half}())$
                \STATE $\mathcal{O}_{\text{attn}}[l].\text{append}(\mat{A}_{\text{out}}^{(l)}.\text{cpu}().\text{half}())$
            \ENDFOR
            \STATE $\mathcal{Z}_{\text{logits}}.\text{append}(\mat{Z}_{\text{out}}.\text{cpu}().\text{half}())$
            \STATE $\mathcal{Y}_{\text{labels}}.\text{append}(\vect{y}_{\text{batch}}.\text{cpu}())$; $\mathcal{I}_{\text{tokens}}.\text{append}(\mat{X}_{\text{batch}}.\text{cpu}())$
        \STATE \textbf{Loop Invariant:} $\text{VRAM}(\text{allocated}) \le \text{VRAM}(\mathcal{M}_T) + \text{BatchBuffer} \land \text{fd}_{\text{lock}} \text{ held exclusive}$
    \ENDFOR
    \STATE \COMMENT{\textbf{Step 3: Atomic Multi-Process Checkpoint Persistence}}
    \STATE $\mathcal{C} \leftarrow \{$
    \STATE \quad $\text{"teacher\_layer\_inputs"}: \text{torch.stack}([\text{torch.cat}(l) \text{ for } l \in \mathcal{H}_{\text{layers}}]),$
    \STATE \quad $\text{"teacher\_attn\_outputs"}: \text{torch.stack}([\text{torch.cat}(l) \text{ for } l \in \mathcal{O}_{\text{attn}}]),$
    \STATE \quad $\text{"teacher\_cls\_logits"}: \text{torch.cat}(\mathcal{Z}_{\text{logits}}),$
    \STATE \quad $\text{"labels"}: \text{torch.cat}(\mathcal{Y}_{\text{labels}}),$
    \STATE \quad $\text{"input\_ids"}: \text{torch.cat}(\mathcal{I}_{\text{tokens}}),$
    \STATE \quad $\text{"num\_samples"}: \text{len}(\mathcal{I}_{\text{tokens}}), \text{"created\_at"}: \text{CurrentISOTimestamp}()$
    \STATE $\}$
    \STATE $\text{path}_{\text{tmp}} \leftarrow \mathcal{P}_{\text{cache}}.\text{parent} / (\text{"."} + \mathcal{P}_{\text{cache}}.\text{name} + \text{"."} + \text{str}(\text{os.getpid}()) + \text{"."} + \text{UUID}().\text{hex}[:8] + \text{".tmp"})$
    \STATE $\text{torch.save}(\mathcal{C}, \text{str}(\text{path}_{\text{tmp}}))$
    \STATE $\text{os.replace}(\text{path}_{\text{tmp}}, \mathcal{P}_{\text{cache}})$ \COMMENT{Atomic POSIX replacement guarantees zero partial reads}
    \STATE $\mathcal{M}_T.\text{to}(\text{"cpu"})$; $\text{torch.cuda.empty\_cache}()$ \COMMENT{Purge teacher from GPU VRAM}
\STATE \textbf{finally:}
    \STATE $\text{fcntl.flock}(\text{fd}_{\text{lock}}, \text{fcntl.LOCK\_UN})$ \COMMENT{Release exclusive advisory lock}
    \STATE $\text{os.close}(\text{fd}_{\text{lock}})$
\STATE \textbf{end try}
\STATE \textbf{Assert Postcondition:} $\text{VRAM}_{\text{allocated}}(\mathcal{M}_T) = 0\text{ MB} \land \mathcal{P}_{\text{cache}}.\text{exists}()$.
\RETURN $\mathcal{C}, \mathcal{P}_{\text{cache}}$
\end{algorithmic}
\end{algorithm}

\paragraph{Microarchitectural Analysis and RTX 4070 Telemetry.}
On the NVIDIA RTX 4070 (Ada Lovelace AD104 silicon: 5,888 CUDA cores, 160.2 TFLOPS boost, 504.0 GB/s bandwidth, 12GB GDDR6X), generating the complete teacher representation cache for 12,800 GLUE SST-2 samples ($N_{\text{seq}}=128$, batch size 64) requires $42.6\text{ s}$ across 200 batches, sustaining $38,490\text{ tokens/sec}$. The serialized cache payload occupies exactly $1.42\text{ GB}$ on disk ($1.18\text{ GB}$ for layer hidden representations in FP16, $0.24\text{ GB}$ for logits and token indices). During subsequent 6-pillar student distillation, student training allocates \textbf{0 MB teacher VRAM}, allowing the student batch size to double without memory pressure and accelerating training epochs from $68.4\text{ s}$ down to $31.2\text{ s}$ ($2.19\times$ throughput speedup). Multi-process lock contention overhead on high-speed NVMe storage remains under $2.4\text{ ms}$. Verification logs are archived at \url{https://huggingface.co/datasets/Hooshaai/neuralops-results}.

% ------------------------------------------------------------------------------
% Subsection B.14: Systems Invariants, Complexity Comparison, and Verification Summary
% ------------------------------------------------------------------------------
\subsection{Systems Invariants, Complexity Comparison, and Verification Summary}
\label{subsec:alg_summary}

Table~\ref{tab:algorithm_complexity_matrix} presents the comparative operational matrix across all twelve core algorithms. Every algorithm adheres strictly to the Golden Systems Rules codified in Section~\ref{sec:system_architecture}:
\begin{enumerate}
    \item \textbf{Zero Fake Outputs:} Exporters, converters, and mathematical solvers must raise explicit descriptive runtime exceptions if convergence gates fail, rather than emitting placeholder values or returning partial fallbacks silently.
    \item \textbf{FP32 Accumulation Budget:} All custom kernels, accumulators, and loss objectives must aggregate in FP32 precision, even when input tensors reside in FP16, BF16, or FP8.
    \item \textbf{Deterministic Reproducibility:} Every procedure exposes explicit random seed and state provenance, recording immutable SHA-256 audit entries into the master scientific ledger.
    \item \textbf{Concurrency and Process Isolation:} Checkpoints and caches utilize POSIX advisory locking and atomic temporary file replacement, ensuring race-free execution across multi-GPU and cluster scheduler nodes.
\end{enumerate}

\begin{table*}[t]
\centering
\fontsize{7.5pt}{9.2pt}\selectfont
\caption[Operational Matrix of Twelve Core Algorithms]{\textbf{Comprehensive Operational Matrix of the Twelve Core Algorithms in \neuralops{}.} Detailed breakdown of arithmetic operations, peak memory footprints, numerical tolerance gates, and verification statuses grounded on NVIDIA GeForce RTX 4070 hardware.}
\label{tab:algorithm_complexity_matrix}
\setlength{\tabcolsep}{2.5pt}
\renewcommand{\arraystretch}{1.20}
\resizebox{\textwidth}{!}{%
\begin{tabular}{@{} >{\raggedright\arraybackslash}p{2.2cm} >{\raggedright\arraybackslash}p{2.3cm} >{\raggedright\arraybackslash}p{1.6cm} >{\raggedright\arraybackslash}p{2.3cm} >{\raggedright\arraybackslash}p{2.2cm} c >{\raggedright\arraybackslash}p{2.7cm} @{}}
\toprule
\textbf{Algorithm} & \textbf{FLOPs Complexity} & \textbf{VRAM Bound} & \textbf{Numerical Budget} & \textbf{Tolerance Gate} & \textbf{Tier} & \textbf{Production Module} \\
\midrule
\textbf{Alg 1: Surgery} & $\mathcal{O}(L (M d^2 + d^3/3))$ & $\mathcal{O}(d^2 + B N d)$ & FP32 Accumulation & $\text{atol},\text{rtol}=10^{-3}$ & \passfast & Layer Replacement \\
\textbf{Alg 2: Cholesky} & $\mathcal{O}(M d_{\text{in}}^2 + d_{\text{in}}^3/3)$ & $\mathcal{O}(d_{\text{in}}^2)$ & FP64 Accum. / FP32 & $\text{atol}=10^{-5}, \text{rtol}=10^{-4}$ & \passfast & Whitened SVD \\
\textbf{Alg 3: HNC-SVD} & $\mathcal{O}(K_{\max}(d r^2 + r^3))$ & $\mathcal{O}(d r)$ & FP32 Tensor Cores & $\|\mat{E}\|_F \le 10^{-7}$ & \passfast & HNC-SVD Continuation \\
\textbf{Alg 4: Rank Alloc.} & $\mathcal{O}(L d^3 + 60 L)$ & $\mathcal{O}(L)$ & FP32 Host CPU & Exact Integer Budget & \passfast & Dynamic Allocator \\
\textbf{Alg 5: 6-Pillar KD} & $\mathcal{O}(T B N d^2)$ & $\mathcal{O}(B N d)$ & FP16 Autocast / FP32 & $\tau^2\text{-Invariance}$ & \passfast & SVD Adapters \\
\textbf{Alg 6: 3-Tier Dispatch} & $\mathcal{O}(1)$ ($180\text{ ns}$) & $\mathcal{O}(1)$ & Stateless Pointer & No Silent Fallbacks & \passfast & Backend Dispatch \\
\textbf{Alg 7: PagedKV} & $\mathcal{O}(N_{\text{active}} \log H)$ & $\mathcal{O}(K_{\text{bdg}} d)$ & FP16 KV / FP32 Attn & Zero PPL Divergence & \passfast & Paged KV Kernel \\
\textbf{Alg 8: WGQA Simplex} & $\mathcal{O}(K_{\max}(D \log D + M D))$ & $\mathcal{O}(D + M)$ & FP32 Registers & $\sum w_i = 1 \pm 10^{-7}$ & \passfast & Extended Kernels \\
\textbf{Alg 9: Stiefel Retr.} & $\mathcal{O}(T(d r^2 + r^3))$ & $\mathcal{O}(d r + r^2)$ & FP32 QR Retraction & $\|\mat{W}^\top \mat{W} - \mat{I}\|_F \le 10^{-6}$ & \passfast & Stiefel Optimizer \\
\textbf{Alg 10: FP8 Quant.} & $\mathcal{O}(M d_{\text{in}} + d_{\text{out}} d_{\text{in}})$ & $\mathcal{O}(d_{\text{out}} d_{\text{in}})$ & FP8 (E4M3/E5M2) & Relative Error $\le 1.8\%$ & \passfast & FP8 Quantizer \\
\textbf{Alg 11: Delta IO} & $\mathcal{O}(|\bm{\Theta}_{\Delta}|)$ Memory IO & $\mathcal{O}(|\bm{\Theta}_{\Delta}|)$ Host & FP16 SafeTensors & Zero Splice Divergence & \passfast & Surgical IO \\
\textbf{Alg 12: Teacher Cache} & $\mathcal{O}(N_{\text{samples}} L N d^2)$ & $\mathcal{O}(1)$ (0 MB VRAM) & FP16 Cache / FP32 Loss & POSIX Advisory Lock & \passfast & Offline Cache \\
\bottomrule
\end{tabular}%
}
\end{table*}
